\documentclass[12pt,a4paper,oneside]{report}

% ==============================================================================
% GÓI VÀ CẤU HÌNH NGÔN NGỮ, FONT CHỮ
% ==============================================================================
\usepackage{iftex}
\ifXeTeX
  \usepackage{fontspec}
  \usepackage{polyglossia}
  \setmainlanguage{vietnamese}
  \setmainfont{Times New Roman}
\else\ifLuaTeX
  \usepackage{fontspec}
  \usepackage{polyglossia}
  \setmainlanguage{vietnamese}
  \setmainfont{Times New Roman}
\else
  \usepackage[T5,T1]{fontenc}
  \usepackage[utf8]{inputenc}
  \usepackage[vietnamese]{babel}
  \usepackage{times}
\fi\fi

% ==============================================================================
% ĐỊNH DẠNG TRANG & CĂN LỀ
% ==============================================================================
\usepackage{geometry}
\geometry{
    a4paper,
    top=25mm,
    bottom=25mm,
    left=30mm,
    right=20mm,
    headheight=15pt
}
\usepackage{setspace}
\onehalfspacing % Giãn dòng 1.5

% ==============================================================================
% GÓI ĐỒ HỌA, TOÁN HỌC, BẢNG BIỂU & MÃ NGUỒN
% ==============================================================================
\usepackage{amsmath, amssymb, amsfonts}
\usepackage{graphicx}
\usepackage{xcolor}
\usepackage{colortbl}
\usepackage{booktabs}
\usepackage{tabularx}
\usepackage{array}
\usepackage{multirow}
\usepackage{longtable}
\usepackage{float}
\usepackage{caption}
\usepackage{subcaption}
\usepackage{enumitem}
\usepackage{listings}
\usepackage[normalem]{ulem}
\usepackage{tikz}
\usetikzlibrary{
    positioning,
    shapes.geometric,
    shapes.misc,
    shapes.multipart,
    arrows.meta,
    calc,
    fit,
    backgrounds,
    shadows.blur,
    shadows
}

% Cấu hình Caption chuyên nghiệp
\captionsetup{
    font=small,
    labelfont={bf, color=navy},
    textfont=it,
    justification=centering,
    skip=8pt
}

% ==============================================================================
% MÀU SẮC & CẤU HÌNH TIKZ / BẢNG BIỂU
% ==============================================================================
\definecolor{navy}{RGB}{15,35,65}
\definecolor{darkblue}{RGB}{20,45,85}
\definecolor{blue}{RGB}{30,90,200}
\definecolor{softblue}{RGB}{238,244,255}
\definecolor{verylight}{RGB}{248,250,253}
\definecolor{graytext}{RGB}{70,80,95}
\definecolor{bordergray}{RGB}{195,208,225}
\definecolor{green}{RGB}{16,150,75}
\definecolor{lightgreen}{RGB}{235,250,240}
\definecolor{amber}{RGB}{215,115,10}
\definecolor{lightamber}{RGB}{254,248,235}
\definecolor{purple}{RGB}{120,50,180}
\definecolor{lightpurple}{RGB}{248,240,255}
\definecolor{tableheader}{RGB}{25,50,90}
\definecolor{tablerowalt}{RGB}{245,248,253}

% Định nghĩa bảng màu học thuật cho ERD
\definecolor{erdnavy}{RGB}{24, 43, 73}
\definecolor{erdheader}{RGB}{238, 244, 250}
\definecolor{erdline}{RGB}{70, 80, 95}
\definecolor{pkcolor}{RGB}{185, 30, 30}
\definecolor{fkcolor}{RGB}{30, 100, 175}

\newcolumntype{L}[1]{>{\raggedright\arraybackslash}p{#1}}
\newcolumntype{C}[1]{>{\centering\arraybackslash}p{#1}}
\newcolumntype{R}[1]{>{\raggedleft\arraybackslash}p{#1}}
\newcolumntype{Y}{>{\raggedright\arraybackslash}X}
\newcolumntype{Z}{>{\centering\arraybackslash}X}

\tikzset{
    erd entity/.style={
        rectangle split,
        rectangle split parts=2,
        rectangle split part fill={erdheader, white},
        rounded corners=2pt,
        draw=erdnavy,
        line width=0.8pt,
        font=\footnotesize\sffamily,
        text width=3.8cm,
        inner sep=4pt,
        rectangle split part align={center, left}
    },
    erd rel/.style={
        draw=erdline,
        line width=0.75pt,
        -{Latex[length=2.2mm]}
    },
    erd card/.style={
        font=\scriptsize\bfseries\sffamily,
        text=erdnavy,
        fill=white,
        inner sep=1.5pt
    },
    entity/.style={
        rectangle,
        rounded corners=4pt,
        minimum width=3.2cm,
        minimum height=1.05cm,
        draw=navy,
        line width=1.3pt,
        fill=softblue,
        align=center,
        font=\bfseries\small,
        text=navy,
        drop shadow={opacity=0.15, shadow xshift=1pt, shadow yshift=-1.5pt}
    },
    weakentity/.style={
        rectangle,
        rounded corners=4pt,
        minimum width=3.2cm,
        minimum height=1.05cm,
        double,
        double distance=1.5pt,
        draw=navy,
        line width=1pt,
        fill=softblue,
        align=center,
        font=\bfseries\small,
        text=navy,
        drop shadow={opacity=0.15, shadow xshift=1pt, shadow yshift=-1.5pt}
    },
    relation/.style={
        diamond,
        aspect=2,
        minimum width=2.2cm,
        minimum height=1.1cm,
        draw=blue,
        line width=1.1pt,
        fill=white,
        align=center,
        font=\scriptsize\bfseries,
        text=navy,
        drop shadow={opacity=0.12, shadow xshift=1pt, shadow yshift=-1pt}
    },
    weakrelation/.style={
        diamond,
        aspect=2,
        minimum width=2.2cm,
        minimum height=1.1cm,
        double,
        double distance=1.2pt,
        draw=blue,
        line width=0.9pt,
        fill=white,
        align=center,
        font=\scriptsize\bfseries,
        text=navy,
        drop shadow={opacity=0.12, shadow xshift=1pt, shadow yshift=-1pt}
    },
    attribute/.style={
        ellipse,
        draw=graytext,
        line width=0.9pt,
        fill=white,
        align=center,
        font=\scriptsize,
        text=navy,
        inner sep=3pt
    },
    pkattribute/.style={
        rectangle,
        rounded corners=3pt,
        draw=green,
        line width=1pt,
        fill=lightgreen,
        align=center,
        font=\scriptsize\bfseries,
        text=green!60!black,
        inner sep=4pt
    },
    connector/.style={
        -{Latex[length=2.2mm, width=1.6mm]},
        line width=1pt,
        draw=graytext!90
    },
    lineconnector/.style={
        line width=1pt,
        draw=graytext!90
    },
    cardinality/.style={
        font=\scriptsize\bfseries,
        fill=white,
        draw=bordergray,
        rounded corners=2pt,
        inner sep=2pt,
        text=blue!80!black
    },
    groupbox/.style={
        rectangle,
        rounded corners=8pt,
        draw=bordergray,
        line width=0.8pt,
        fill=verylight,
        inner sep=10pt
    }
}

% Định dạng khối mã nguồn (SQL, Python, v.v.)
\lstset{
    basicstyle=\ttfamily\small,
    keywordstyle=\color{blue}\bfseries,
    stringstyle=\color{red!70!black},
    commentstyle=\color{green!50!black}\itshape,
    numbers=left,
    numberstyle=\tiny\color{gray},
    stepnumber=1,
    numbersep=8pt,
    backgroundcolor=\color{gray!5},
    showspaces=false,
    showstringspaces=false,
    showtabs=false,
    frame=single,
    rulecolor=\color{gray!30},
    tabsize=2,
    captionpos=b,
    breaklines=true,
    breakatwhitespace=false
}

% ==============================================================================
% HEADER & FOOTER
% ==============================================================================
\usepackage{fancyhdr}
\pagestyle{fancy}
\fancyhf{}
\fancyhead[L]{\small\slshape Đồ án Thiết kế Cơ sở Dữ liệu}
\fancyhead[R]{\small\slshape Đề tài: FamilyConnect}
\fancyfoot[C]{\thepage}
\renewcommand{\headrulewidth}{0.4pt}
\renewcommand{\footrulewidth}{0.4pt}

% Áp dụng plain style cho trang đầu chương
\fancypagestyle{plain}{
  \fancyhf{}
  \fancyfoot[C]{\thepage}
  \renewcommand{\headrulewidth}{0pt}
  \renewcommand{\footrulewidth}{0pt}
}

% ==============================================================================
% TIÊU ĐỀ CHƯƠNG & MỤC
% ==============================================================================
\usepackage{titlesec}
\titleformat{\chapter}[hang]
  {\normalfont\LARGE\bfseries}{\chaptertitlename\ \thechapter:}{0.8em}{}
\titlespacing*{\chapter}{0pt}{-20pt}{20pt}

% ==============================================================================
% HYPERREF (LIÊN KẾT & MỤC LỤC TƯƠNG TÁC)
% ==============================================================================
\usepackage{hyperref}
\hypersetup{
    colorlinks=true,
    linkcolor=black,
    filecolor=blue,
    citecolor=blue,
    urlcolor=blue,
    pdftitle={Đồ án Thiết kế Cơ sở Dữ liệu - FamilyConnect},
    pdfauthor={Nhóm sinh viên},
    pdfsubject={Database Design Project},
    pdfkeywords={Database, ERD, Relational Model, Normalization, SQL}
}

% ==============================================================================
% BẮT ĐẦU TÀI LIỆU
% ==============================================================================
\begin{document}

% ------------------------------------------------------------------------------
% TRANG BÌA
% ------------------------------------------------------------------------------
\hypersetup{pageanchor=false}
\begin{titlepage}
\begin{center}
    \large\textbf{TRƯỜNG ĐẠI HỌC GIAO THÔNG VẬN TẢI TP. HỒ CHÍ MINH}\\
    \textbf{VIỆN ĐÀO TẠO / KHOA CÔNG NGHỆ THÔNG TIN}\\[1.5cm]

    % Logo (nếu có ảnh logo, có thể bỏ comment dòng dưới)
    % \includegraphics[width=0.25\textwidth]{images/logo_uth.png}\\[1cm]
    \vspace*{1.5cm}

    \Large\textbf{BÁO CÁO ĐỒ ÁN MÔN HỌC}\\[0.5cm]
    \Huge\textbf{THIẾT KẾ CƠ SỞ DỮ LIỆU}\\[1.5cm]

    \rule{\linewidth}{1.5pt}\\[0.4cm]
    \LARGE\textbf{ĐỀ TÀI: FAMILYCONNECT}\\[0.2cm]
    \large\textbf{Nền tảng Cộng đồng Gia đình số tích hợp Trí tuệ Nhân tạo}\\
    \rule{\linewidth}{1.5pt}\\[2cm]

    \begin{minipage}[t]{0.45\textwidth}
        \textbf{Giảng viên hướng dẫn:}\\
        \indent ThS. / TS. [Tên Giảng Viên]
    \end{minipage}
    \begin{minipage}[t]{0.5\textwidth}
        \textbf{Sinh viên thực hiện:}\\
        \indent 1. Trương Minh Nhựt\\
        \indent 2. Phan Thủy Tiên\\
        \indent 3. Nguyễn Tuấn Kiệt\\
        \indent 4. Trần Quốc Khánh\\
        \indent 5. Huỳnh Như
    \end{minipage}

    \vfill
    {\large TP. Hồ Chí Minh, Năm 2026}
\end{center}
\end{titlepage}
\hypersetup{pageanchor=true}

% ------------------------------------------------------------------------------
% MỤC LỤC VÀ DANH MỤC
% ------------------------------------------------------------------------------
\pagenumbering{roman}
\tableofcontents
\newpage
\listoffigures
\newpage
\listoftables
\newpage

% ------------------------------------------------------------------------------
% PHẦN NỘI DUNG CHÍNH (CÁC CHƯƠNG)
% ------------------------------------------------------------------------------
\pagenumbering{arabic}

% Giới thiệu đề tài (Không đánh số chương)
% !TEX root = ../main.tex

\chapter*{GIỚI THIỆU ĐỀ TÀI}

\addcontentsline{toc}{chapter}{GIỚI THIỆU ĐỀ TÀI}

\label{chap:gioi_thieu_de_tai}

% =========================================================
% 1. BỐI CẢNH VÀ LÝ DO CHỌN ĐỀ TÀI
% =========================================================

\section*{1. Bối cảnh và lý do chọn đề tài}

\addcontentsline{toc}{section}{1. Bối cảnh và lý do chọn đề tài}

Trong bối cảnh xã hội hiện đại, sự phát triển mạnh mẽ của công nghệ thông tin và quá trình toàn cầu hóa đã tạo ra nhiều thay đổi trong đời sống gia đình và cộng đồng. Các thành viên trong gia đình ngày càng có xu hướng sinh sống, học tập và làm việc tại nhiều tỉnh thành khác nhau, thậm chí ở các quốc gia khác nhau. Điều này mang lại nhiều cơ hội phát triển cá nhân nhưng đồng thời cũng đặt ra thách thức trong việc duy trì sự gắn kết giữa các thế hệ, lưu giữ truyền thống gia đình và bảo tồn thông tin về nguồn gốc dòng họ.

Tại Việt Nam, gia đình và dòng họ giữ vai trò quan trọng trong việc duy trì các giá trị văn hóa, đạo đức và truyền thống dân tộc. Những câu nói như ``Uống nước nhớ nguồn'' hay ``Chim có tổ, người có tông'' thể hiện rõ ý nghĩa của việc ghi nhớ cội nguồn và duy trì mối liên kết giữa các thế hệ. Tuy nhiên, việc quản lý thông tin gia đình, gia phả và các hoạt động dòng họ hiện nay vẫn còn phụ thuộc nhiều vào phương pháp truyền thống hoặc các công cụ chưa được thiết kế chuyên biệt.

Qua khảo sát thực tế, các giải pháp hiện có còn tồn tại một số hạn chế. Các mạng xã hội phổ biến như Facebook, Zalo hay TikTok chủ yếu phục vụ nhu cầu giao tiếp chung, chưa hỗ trợ tốt việc tổ chức và lưu trữ dữ liệu có cấu trúc về gia đình, dòng họ và các mối quan hệ huyết thống. Bên cạnh đó, các phần mềm gia phả truyền thống thường tập trung vào việc tra cứu thông tin tĩnh, thiếu các tính năng tương tác cộng đồng, tổ chức sự kiện và chia sẻ di sản gia đình.

Từ những thực trạng trên, việc xây dựng một nền tảng số chuyên biệt nhằm kết nối các thành viên trong gia đình, quản lý thông tin gia phả và lưu giữ các giá trị truyền thống là cần thiết. Đề tài FamilyConnect được đề xuất nhằm giải quyết những hạn chế trên thông qua việc kết hợp giữa quản lý dữ liệu gia đình, tương tác cộng đồng và ứng dụng công nghệ trí tuệ nhân tạo.

Đặc biệt, đề tài tập trung nghiên cứu và xây dựng cơ sở dữ liệu làm nền tảng cho toàn bộ hệ thống. Cơ sở dữ liệu cần đảm bảo khả năng lưu trữ chính xác thông tin người dùng, gia đình, thành viên, chi nhánh dòng họ, bài đăng, sự kiện và di sản. Đồng thời, hệ thống phải duy trì được tính toàn vẹn dữ liệu và phản ánh chính xác các mối quan hệ giữa các thực thể trong môi trường gia đình và dòng tộc.

% =========================================================
% 2. MỤC TIÊU VÀ GIẢI PHÁP
% =========================================================

\section*{2. Mục tiêu đề tài \& Giải pháp đề xuất}

\addcontentsline{toc}{section}{2. Mục tiêu đề tài \& Giải pháp đề xuất}

\subsection*{2.1. Mục tiêu đề tài}

\addcontentsline{toc}{subsection}{2.1. Mục tiêu đề tài}

Mục tiêu tổng quát của đề tài là nghiên cứu, phân tích và thiết kế một hệ thống cơ sở dữ liệu phục vụ nền tảng FamilyConnect, nhằm hỗ trợ việc quản lý thông tin gia đình, dòng họ và các hoạt động cộng đồng trên môi trường số.

Đề tài hướng đến các mục tiêu cụ thể sau:

\begin{itemize}[leftmargin=*]
\item Xây dựng mô hình dữ liệu phản ánh đầy đủ các đối tượng nghiệp vụ trong hệ thống FamilyConnect.
\item Xác định các thực thể, thuộc tính và mối quan hệ giữa các đối tượng dữ liệu.
\item Thiết kế cơ sở dữ liệu có cấu trúc rõ ràng, hạn chế dư thừa và đảm bảo tính nhất quán của dữ liệu.
\item Quản lý hiệu quả thông tin thành viên, gia đình, chi nhánh dòng họ và các mối quan hệ huyết thống.
\item Hỗ trợ lưu trữ và quản lý các bài đăng, sự kiện và di sản gia đình.
\item Đảm bảo khả năng mở rộng cơ sở dữ liệu để phục vụ các chức năng phát triển trong tương lai, bao gồm tích hợp trí tuệ nhân tạo.
\end{itemize}

\subsection*{2.2. Giải pháp đề xuất}

\addcontentsline{toc}{subsection}{2.2. Giải pháp đề xuất}

FamilyConnect được đề xuất là một nền tảng cộng đồng gia đình số kết hợp giữa quản lý dữ liệu gia phả và các chức năng tương tác hiện đại. Giải pháp được xây dựng dựa trên việc thiết kế một hệ thống cơ sở dữ liệu tập trung, đóng vai trò lưu trữ và cung cấp thông tin cho các chức năng của ứng dụng.

Các nhóm chức năng chính của hệ thống bao gồm:

\begin{enumerate}[leftmargin=*]
\item \textbf{Quản lý người dùng và gia đình:} Lưu trữ thông tin tài khoản, hồ sơ cá nhân và thông tin các gia đình tham gia hệ thống.
\item \textbf{Quản lý thành viên và gia phả:} Quản lý thông tin thành viên, các mối quan hệ giữa các thành viên và cấu trúc dòng họ qua nhiều thế hệ.
\item \textbf{Quản lý chi nhánh dòng họ:} Hỗ trợ tổ chức thông tin theo các chi, nhánh hoặc đơn vị gia đình thuộc cùng một dòng họ.
\item \textbf{Quản lý bài đăng và tương tác cộng đồng:} Cho phép thành viên chia sẻ thông tin, hình ảnh, câu chuyện và hoạt động trong không gian riêng tư của gia đình.
\item \textbf{Quản lý sự kiện:} Lưu trữ và tổ chức các sự kiện gia đình, dòng họ như họp mặt, ngày giỗ, lễ tết hoặc các hoạt động cộng đồng.
\item \textbf{Quản lý di sản gia đình:} Lưu trữ các hình ảnh, tài liệu, câu chuyện và tư liệu có giá trị lịch sử, văn hóa của gia đình và dòng họ.
\item \textbf{Tích hợp trí tuệ nhân tạo:} Hỗ trợ tìm kiếm thông tin ngữ nghĩa và giải thích các mối quan hệ dòng tộc dựa trên dữ liệu được lưu trữ trong hệ thống.
\end{enumerate}

Giải pháp đề xuất không chỉ hướng đến việc số hóa thông tin gia phả mà còn xây dựng một môi trường kết nối cộng đồng gia đình. Trong đó, cơ sở dữ liệu giữ vai trò trung tâm, đảm bảo dữ liệu được tổ chức khoa học, chính xác và có khả năng phục vụ nhiều nhóm chức năng khác nhau.

% =========================================================
% 3. PHẠM VI VÀ ĐỐI TƯỢNG
% =========================================================

\section*{3. Phạm vi và đối tượng áp dụng}

\addcontentsline{toc}{section}{3. Phạm vi và đối tượng áp dụng}

\subsection*{3.1. Đối tượng áp dụng}

\addcontentsline{toc}{subsection}{3.1. Đối tượng áp dụng}

Hệ thống FamilyConnect được hướng đến các nhóm đối tượng sử dụng sau:

\begin{itemize}[leftmargin=*]
\item \textbf{Thành viên gia đình:} Các cá nhân thuộc gia đình hạt nhân hoặc gia đình mở rộng có nhu cầu kết nối, chia sẻ thông tin và duy trì quan hệ giữa các thế hệ.
\item \textbf{Dòng họ và chi nhánh gia tộc:} Các dòng họ, chi họ hoặc nhánh gia đình có nhu cầu quản lý thông tin thành viên, cây gia phả và các hoạt động chung.
\item \textbf{Ban quản lý dòng họ:} Trưởng họ, người quản lý gia phả hoặc những cá nhân được phân quyền quản trị thông tin và tổ chức các hoạt động cộng đồng.
\item \textbf{Người quản trị hệ thống:} Người chịu trách nhiệm quản lý tài khoản, phân quyền và đảm bảo hoạt động ổn định của hệ thống.
\end{itemize}

\subsection*{3.2. Phạm vi đề tài}

\addcontentsline{toc}{subsection}{3.2. Phạm vi đề tài}

Trong khuôn khổ môn học Thiết kế Cơ sở dữ liệu, đề tài tập trung chủ yếu vào việc khảo sát, phân tích và thiết kế cơ sở dữ liệu cho hệ thống FamilyConnect.

Phạm vi nghiên cứu bao gồm:

\begin{itemize}[leftmargin=*]
\item Phân tích yêu cầu nghiệp vụ của nền tảng FamilyConnect.
\item Xác định các thực thể và thuộc tính cần quản lý trong hệ thống.
\item Xây dựng mô hình thực thể kết hợp (ER Model).
\item Xác định các mối quan hệ giữa các thực thể và các ràng buộc dữ liệu.
\item Chuyển đổi mô hình thực thể kết hợp sang mô hình quan hệ.
\item Thiết kế các bảng dữ liệu, khóa chính, khóa ngoại và các ràng buộc toàn vẹn.
\item Chuẩn hóa cơ sở dữ liệu nhằm hạn chế dư thừa dữ liệu và đảm bảo tính nhất quán.
\item Định hướng triển khai cơ sở dữ liệu trên hệ quản trị cơ sở dữ liệu MySQL.
\end{itemize}

Đề tài không đi sâu vào việc phát triển hoàn chỉnh giao diện người dùng hoặc triển khai các mô hình trí tuệ nhân tạo. Các chức năng AI được xem xét ở mức độ định hướng tích hợp và khai thác dữ liệu từ cơ sở dữ liệu đã được thiết kế.

% =========================================================
% 4. KHẢO SÁT HIỆN TRẠNG VÀ YÊU CẦU HỆ THỐNG
% =========================================================

\section*{4. Khảo sát hiện trạng \& Các yêu cầu của hệ thống}

\addcontentsline{toc}{section}{4. Khảo sát hiện trạng \& Các yêu cầu của hệ thống}

\subsection*{4.1. Yêu cầu chức năng}

\addcontentsline{toc}{subsection}{4.1. Yêu cầu chức năng}

Dựa trên quá trình khảo sát nhu cầu sử dụng và mục tiêu của hệ thống, FamilyConnect cần đáp ứng các yêu cầu chức năng chính sau:

\begin{enumerate}[leftmargin=*]
\item \textbf{Quản lý người dùng:} Cho phép đăng ký, đăng nhập, cập nhật và quản lý thông tin cá nhân của người dùng.
\item \textbf{Quản lý gia đình:} Cho phép tạo, cập nhật và quản lý thông tin các gia đình hoặc dòng họ tham gia hệ thống.
\item \textbf{Quản lý thành viên:} Cho phép thêm, sửa, xóa và tra cứu thông tin thành viên thuộc một gia đình hoặc dòng họ.
\item \textbf{Quản lý quan hệ gia phả:} Cho phép xác định và lưu trữ các mối quan hệ giữa các thành viên như cha mẹ, con cái, vợ chồng hoặc các quan hệ huyết thống khác.
\item \textbf{Quản lý chi nhánh dòng họ:} Cho phép phân chia và quản lý các chi, nhánh thuộc một dòng họ.
\item \textbf{Quản lý bài đăng:} Cho phép thành viên tạo, chỉnh sửa, xóa và chia sẻ bài viết trong phạm vi gia đình hoặc dòng họ.
\item \textbf{Quản lý tương tác:} Cho phép người dùng bình luận và tương tác với các bài đăng được chia sẻ trong cộng đồng.
\item \textbf{Quản lý sự kiện:} Cho phép tạo, cập nhật và theo dõi các sự kiện gia đình, đồng thời quản lý danh sách thành viên tham dự.
\item \textbf{Quản lý di sản:} Cho phép lưu trữ và tra cứu các tài liệu, hình ảnh, câu chuyện và thông tin lịch sử liên quan đến gia đình.
\item \textbf{Quản lý phân quyền:} Cho phép phân quyền người dùng theo các vai trò khác nhau như quản trị viên, trưởng họ, thành viên và khách.
\item \textbf{Tra cứu và tìm kiếm thông tin:} Hỗ trợ tìm kiếm thành viên, gia đình, sự kiện và các thông tin liên quan trong hệ thống.
\item \textbf{Hỗ trợ AI:} Cho phép khai thác dữ liệu để phục vụ chức năng tìm kiếm ngữ nghĩa và giải thích mối quan hệ giữa các thành viên.
\end{enumerate}

Các yêu cầu chức năng trên là cơ sở để xác định các thực thể, thuộc tính và mối quan hệ cần thiết trong quá trình thiết kế cơ sở dữ liệu.

\subsection*{4.2. Yêu cầu phi chức năng}

\addcontentsline{toc}{subsection}{4.2. Yêu cầu phi chức năng}

Bên cạnh các yêu cầu chức năng, hệ thống FamilyConnect cần đáp ứng các yêu cầu phi chức năng nhằm đảm bảo chất lượng, tính ổn định và khả năng phát triển lâu dài.

\begin{itemize}[leftmargin=*]
\item \textbf{Tính toàn vẹn dữ liệu:} Cơ sở dữ liệu phải đảm bảo tính chính xác, nhất quán và hợp lệ của thông tin. Các ràng buộc về khóa chính, khóa ngoại và quan hệ giữa các thực thể cần được kiểm soát chặt chẽ.
\item \textbf{Tính bảo mật:} Thông tin cá nhân, dữ liệu gia đình và thông tin dòng họ cần được bảo vệ khỏi việc truy cập trái phép. Hệ thống phải hỗ trợ cơ chế phân quyền phù hợp với từng nhóm người dùng.
\item \textbf{Hiệu năng:} Hệ thống cần đảm bảo thời gian truy vấn và phản hồi hợp lý khi số lượng thành viên, gia đình và dữ liệu gia phả tăng lên.
\item \textbf{Khả năng mở rộng:} Thiết kế cơ sở dữ liệu cần có khả năng mở rộng để bổ sung các chức năng mới trong tương lai mà không làm ảnh hưởng đến cấu trúc dữ liệu hiện tại.
\item \textbf{Tính nhất quán dữ liệu:} Các thông tin liên quan giữa người dùng, thành viên, gia đình và các hoạt động cộng đồng phải được duy trì đồng bộ trong toàn hệ thống.
\item \textbf{Khả năng bảo trì:} Cấu trúc cơ sở dữ liệu cần được tổ chức rõ ràng, dễ hiểu và thuận tiện cho việc bảo trì, cập nhật và phát triển hệ thống.
\item \textbf{Khả năng sao lưu và phục hồi:} Dữ liệu cần có cơ chế sao lưu và phục hồi nhằm hạn chế nguy cơ mất mát thông tin quan trọng.
\end{itemize}

Từ các yêu cầu trên, việc thiết kế cơ sở dữ liệu FamilyConnect cần đảm bảo sự cân bằng giữa tính chính xác của dữ liệu, hiệu năng truy vấn và khả năng mở rộng trong tương lai. Đây sẽ là nền tảng quan trọng cho các chương tiếp theo về thiết kế quan niệm, thiết kế logic và triển khai cơ sở dữ liệu.


% Chương 1: Thiết kế Quan niệm (Conceptual Design)
% !TEX root = ../../main.tex
\chapter{Thiết kế Quan niệm (Conceptual Design)}
\label{chap:chuong1_quan_niem}

Thiết kế quan niệm là bước đầu tiên và quan trọng nhất trong quy trình thiết kế cơ sở dữ liệu cho nền tảng \textbf{FamilyConnect}. Mục tiêu của giai đoạn này là mô hình hóa toàn bộ thế giới thực của nghiệp vụ quản lý gia phả, kết nối cộng đồng dòng họ và bảo tồn di sản số thành các cấu trúc dữ liệu trừu tượng độc lập với hệ quản trị cơ sở dữ liệu (DBMS).

% !TEX root = ../../main.tex
\section{Xác định các thực thể và tập thuộc tính (Conceptual Design)}
\label{sec:1_1_thuc_the_thuoc_tinh}

Trong giai đoạn thiết kế quan niệm cho nền tảng \textbf{FamilyConnect}, dựa trên yêu cầu chức năng về quản lý phân quyền, cấu trúc cây gia phả, hoạt động cộng đồng và lưu trữ di sản số, hệ thống được phân rã thành 4 nhóm thực thể chính gồm 15 thực thể nghiệp vụ cốt lõi.

\begin{table}[H]
    \centering
    \small
    \begin{tabularx}{\textwidth}{c l X}
        \toprule
        \rowcolor{tableheader}\textcolor{white}{\textbf{Ký hiệu}} & \textcolor{white}{\textbf{Phân loại thuộc tính}} & \textcolor{white}{\textbf{Ý nghĩa và Quy tắc nghiệp vụ}} \\
        \midrule
        \rowcolor{tablerowalt}
        \textbf{($*$)} & Thuộc tính bắt buộc & Giá trị bắt buộc phải có khi khởi tạo đối tượng (\texttt{Not Null}). \\
        \textbf{($\circ$)} & Thuộc tính tùy chọn & Giá trị không bắt buộc, có thể bổ sung khi có thông tin (\texttt{Optional}). \\
        \rowcolor{tablerowalt}
        \textbf{(Derived)} & Thuộc tính suy diễn & Được tính toán tự động từ thuộc tính khác (ví dụ: \texttt{Tuoi} tính từ \texttt{NgaySinh}). \\
        \bottomrule
    \end{tabularx}
    \caption{Quy ước ký hiệu thuộc tính trong thiết kế quan niệm}
    \label{tab:attribute_notations}
\end{table}

% ==============================================================================
% 1. NHÓM NGƯỜI DÙNG & BẢO MẬT
% ==============================================================================
\subsection{Nhóm Thực thể Người dùng \& Bảo mật (User \& Security)}
Nhóm thực thể này đảm bảo việc xác thực tài khoản, kiểm soát truy cập dựa trên vai trò (Role-Based Access Control -- RBAC) và ghi nhận nhật ký vận hành nhằm bảo đảm an toàn hệ thống.

\subsubsection{Thực thể NGUOI\_DUNG (User)}
Đại diện cho tài khoản đăng nhập vào hệ thống (thông qua Web hoặc Mobile App). Một tài khoản có thể liên kết với một hồ sơ thành viên trong gia tộc sau khi được xác thực.

\begin{table}[H]
    \centering
    \small
    \begin{tabularx}{\textwidth}{l c X p{4.2cm}}
        \toprule
        \rowcolor{tableheader}\textcolor{white}{\textbf{Tên thuộc tính}} & \textcolor{white}{\textbf{Ký hiệu}} & \textcolor{white}{\textbf{Mô tả ý nghĩa nghiệp vụ}} & \textcolor{white}{\textbf{Ví dụ minh họa}} \\
        \midrule
        \rowcolor{tablerowalt}
        \texttt{MaNguoiDung} & $*$ & Mã định danh duy nhất của người dùng & \texttt{USR\_001}, \texttt{USR\_092} \\
        \texttt{Email} & $*$ & Địa chỉ thư điện tử dùng để đăng nhập và nhận thông báo & \texttt{tuan.nguyen@gmail.com} \\
        \rowcolor{tablerowalt}
        \texttt{SoDienThoai} & $*$ & Số điện thoại đăng ký/xác thực tài khoản & \texttt{0901234567} \\
        \texttt{MatKhauHash} & $*$ & Mật khẩu đã được mã hóa an toàn & \texttt{\$2a\$12\$e8x9A...} \\
        \rowcolor{tablerowalt}
        \texttt{TrangThai} & $*$ & Trạng thái tài khoản người dùng & Đang hoạt động, Chờ duyệt \\
        \texttt{LanDangNhapCuoi} & $\circ$ & Thời điểm đăng nhập gần nhất vào hệ thống & 2026-09-08 08:30:00 \\
        \rowcolor{tablerowalt}
        \texttt{NgayTao} & $*$ & Thời gian khởi tạo tài khoản & 2026-01-15 10:00:00 \\
        \bottomrule
    \end{tabularx}
    \caption{Danh mục thuộc tính thực thể NGUOI\_DUNG (User)}
    \label{tab:ent_nguoi_dung}
\end{table}

\subsubsection{Thực thể VAI\_TRO (Role)}
Xác định nhóm vai trò của người dùng trong hệ thống (ví dụ: Quản trị viên hệ thống -- Super Admin, Trưởng ban liên lạc dòng họ -- Clan Admin, Trưởng chi -- Branch Admin, Thành viên thông thường -- Member).

\begin{table}[H]
    \centering
    \small
    \begin{tabularx}{\textwidth}{l c X p{4.2cm}}
        \toprule
        \rowcolor{tableheader}\textcolor{white}{\textbf{Tên thuộc tính}} & \textcolor{white}{\textbf{Ký hiệu}} & \textcolor{white}{\textbf{Mô tả ý nghĩa nghiệp vụ}} & \textcolor{white}{\textbf{Ví dụ minh họa}} \\
        \midrule
        \rowcolor{tablerowalt}
        \texttt{MaVaiTro} & $*$ & Ký hiệu mã vai trò trong hệ thống & \texttt{ROLE\_CLAN\_ADMIN}, \texttt{ROLE\_MEMBER} \\
        \texttt{TenVaiTro} & $*$ & Tên hiển thị của vai trò & Trưởng ban liên lạc dòng họ, Thành viên \\
        \rowcolor{tablerowalt}
        \texttt{MoTa} & $\circ$ & Giải thích quyền hạn và trách nhiệm của vai trò & Quản trị toàn bộ cây gia phả và duyệt bài viết \\
        \bottomrule
    \end{tabularx}
    \caption{Danh mục thuộc tính thực thể VAI\_TRO (Role)}
    \label{tab:ent_vai_tro}
\end{table}

\subsubsection{Thực thể QUYEN (Permission)}
Quy định chi tiết các hành động nghiệp vụ cụ thể được phép thực thi trên hệ thống (ví dụ: \texttt{genealogy.edit}, \texttt{post.approve}, \texttt{event.create}, \texttt{heritage.delete}).

\begin{table}[H]
    \centering
    \small
    \begin{tabularx}{\textwidth}{l c X p{4.2cm}}
        \toprule
        \rowcolor{tableheader}\textcolor{white}{\textbf{Tên thuộc tính}} & \textcolor{white}{\textbf{Ký hiệu}} & \textcolor{white}{\textbf{Mô tả ý nghĩa nghiệp vụ}} & \textcolor{white}{\textbf{Ví dụ minh họa}} \\
        \midrule
        \rowcolor{tablerowalt}
        \texttt{MaQuyen} & $*$ & Ký hiệu mã quyền duy nhất & \texttt{genealogy.edit}, \texttt{post.approve} \\
        \texttt{TenQuyen} & $*$ & Tên hành động hoặc quyền hạn chi tiết & Chỉnh sửa cây phả hệ, Duyệt bài viết \\
        \rowcolor{tablerowalt}
        \texttt{NhomChucNang} & $*$ & Phân hệ chức năng áp dụng & Gia phả, Bài viết, Sự kiện, Di sản \\
        \bottomrule
    \end{tabularx}
    \caption{Danh mục thuộc tính thực thể QUYEN (Permission)}
    \label{tab:ent_quyen}
\end{table}

\subsubsection{Thực thể NHAT\_KY\_HE\_THONG (Audit Log)}
Lưu lại lịch sử các thao tác nhạy cảm (thêm/sửa/xóa phả hệ, phân quyền, duyệt bài viết) để phục vụ kiểm toán, đối soát và khôi phục khi có sự cố.

\begin{table}[H]
    \centering
    \small
    \begin{tabularx}{\textwidth}{l c X p{4.2cm}}
        \toprule
        \rowcolor{tableheader}\textcolor{white}{\textbf{Tên thuộc tính}} & \textcolor{white}{\textbf{Ký hiệu}} & \textcolor{white}{\textbf{Mô tả ý nghĩa nghiệp vụ}} & \textcolor{white}{\textbf{Ví dụ minh họa}} \\
        \midrule
        \rowcolor{tablerowalt}
        \texttt{MaNhatKy} & $*$ & Mã bản ghi nhật ký vận hành & \texttt{LOG\_100293} \\
        \texttt{ThoiDiem} & $*$ & Thời gian thao tác diễn ra & 2026-09-08 09:15:20 \\
        \rowcolor{tablerowalt}
        \texttt{HanhDong} & $*$ & Thao tác được thực hiện & \texttt{CREATE}, \texttt{UPDATE}, \texttt{DELETE} \\
        \texttt{DoiTuongTacDong} & $*$ & Tên đối tượng hoặc thực thể bị tác động & \texttt{THANH\_VIEN}, \texttt{BAI\_VIET} \\
        \rowcolor{tablerowalt}
        \texttt{DuLieuCu} & $\circ$ & Trạng thái dữ liệu trước khi thay đổi & \texttt{\{"TinhTrang": "Còn sống"\}} \\
        \texttt{DuLieuMoi} & $\circ$ & Trạng thái dữ liệu sau khi cập nhật & \texttt{\{"TinhTrang": "Đã mất"\}} \\
        \rowcolor{tablerowalt}
        \texttt{DiaChiIP} & $\circ$ & Địa chỉ mạng của người dùng thao tác & \texttt{192.168.1.50} \\
        \bottomrule
    \end{tabularx}
    \caption{Danh mục thuộc tính thực thể NHAT\_KY\_HE\_THONG (Audit Log)}
    \label{tab:ent_nhat_ky}
\end{table}

% ==============================================================================
% 2. NHÓM GIA TỘC & GIA PHẢ
% ==============================================================================
\subsection{Nhóm Thực thể Gia tộc \& Gia phả (Family \& Genealogy)}
Đây là nhóm thực thể cốt lõi biểu diễn cây phả hệ, thứ bậc, hồ sơ cá nhân và thông tin nghề nghiệp/học vấn của các thế hệ trong dòng họ.

\subsubsection{Thực thể DONG\_HO (Family / Clan)}
Lưu trữ thông tin khái quát về một đại dòng họ hoặc tổ chức gia tộc lớn.

\begin{table}[H]
    \centering
    \small
    \begin{tabularx}{\textwidth}{l c X p{4.2cm}}
        \toprule
        \rowcolor{tableheader}\textcolor{white}{\textbf{Tên thuộc tính}} & \textcolor{white}{\textbf{Ký hiệu}} & \textcolor{white}{\textbf{Mô tả ý nghĩa nghiệp vụ}} & \textcolor{white}{\textbf{Ví dụ minh họa}} \\
        \midrule
        \rowcolor{tablerowalt}
        \texttt{MaDongHo} & $*$ & Mã định danh dòng họ & \texttt{CLAN\_NGUYEN\_CANH} \\
        \texttt{TenDongHo} & $*$ & Tên dòng họ đại gia tộc & "Họ Nguyễn Cảnh", "Họ Trần Đại Tôn" \\
        \rowcolor{tablerowalt}
        \texttt{ThuyTo} & $\circ$ & Tên cụ Thủy tổ / Khởi tổ của dòng họ & Cụ Nguyễn Cảnh Chân \\
        \texttt{QueQuan} & $*$ & Quê gốc / Nơi phát tích của dòng họ & Làng Đô Lương, Nghệ An \\
        \rowcolor{tablerowalt}
        \texttt{NhaThoTo} & $\circ$ & Địa chỉ hoặc tọa độ từ đường chính & Xã Tràng Sơn, Huyện Đô Lương \\
        \texttt{NgayGioTo} & $\circ$ & Ngày giỗ tổ hàng năm (Âm lịch) & Ngày 15 tháng Giêng \\
        \rowcolor{tablerowalt}
        \texttt{LichSuHinhThanh} & $\circ$ & Tóm tắt lịch sử hình thành và phát triển & Khởi phát từ thế kỷ XV, trải qua 18 đời \\
        \bottomrule
    \end{tabularx}
    \caption{Danh mục thuộc tính thực thể DONG\_HO (Family / Clan)}
    \label{tab:ent_dong_ho}
\end{table}

\subsubsection{Thực thể CHI\_TOC (Branch)}
Biểu diễn các chi, nhánh, phái được phân tách từ dòng họ chính qua các đời.

\begin{table}[H]
    \centering
    \small
    \begin{tabularx}{\textwidth}{l c X p{4.2cm}}
        \toprule
        \rowcolor{tableheader}\textcolor{white}{\textbf{Tên thuộc tính}} & \textcolor{white}{\textbf{Ký hiệu}} & \textcolor{white}{\textbf{Mô tả ý nghĩa nghiệp vụ}} & \textcolor{white}{\textbf{Ví dụ minh họa}} \\
        \midrule
        \rowcolor{tablerowalt}
        \texttt{MaChiToc} & $*$ & Mã định danh chi/nhánh & \texttt{BRANCH\_01\_TRUONG} \\
        \texttt{TenChiToc} & $*$ & Tên gọi của chi họ & "Chi Trưởng", "Chi Hai -- Nhánh Cụ Ba" \\
        \rowcolor{tablerowalt}
        \texttt{DoiThu} & $*$ & Bắt đầu phân tách từ đời thứ mấy & Đời thứ 5 \\
        \texttt{DiaBanChinh} & $\circ$ & Địa bàn sinh sống tập trung hiện tại & TP. Hồ Chí Minh, Hà Nội \\
        \rowcolor{tablerowalt}
        \texttt{GhiChu} & $\circ$ & Thông tin ghi chú thêm về chi nhánh & Chi di cư vào Nam từ năm 1954 \\
        \bottomrule
    \end{tabularx}
    \caption{Danh mục thuộc tính thực thể CHI\_TOC (Branch)}
    \label{tab:ent_chi_toc}
\end{table}

\subsubsection{Thực thể THANH\_VIEN (Member / Profile)}
Mô tả một cá nhân cụ thể trong phả hệ (dù còn sống hay đã mất, dù đã có tài khoản \texttt{NGUOI\_DUNG} hay chỉ là bản ghi lịch sử được ghi chép lại).

\begin{table}[H]
    \centering
    \small
    \begin{tabularx}{\textwidth}{l c X p{4.2cm}}
        \toprule
        \rowcolor{tableheader}\textcolor{white}{\textbf{Tên thuộc tính}} & \textcolor{white}{\textbf{Ký hiệu}} & \textcolor{white}{\textbf{Mô tả ý nghĩa nghiệp vụ}} & \textcolor{white}{\textbf{Ví dụ minh họa}} \\
        \midrule
        \rowcolor{tablerowalt}
        \texttt{MaThanhVien} & $*$ & Mã định danh thành viên & \texttt{MEM\_001}, \texttt{MEM\_092} \\
        \texttt{HoVaTen} & $*$ & Họ và tên khai sinh đầy đủ & Nguyễn Cảnh Tuấn \\
        \rowcolor{tablerowalt}
        \texttt{TenThuongGoi} & $\circ$ & Tên chữ, biệt hiệu, tên gọi ở nhà & Tự Đức Minh, Bé Bo \\
        \texttt{GioiTinh} & $*$ & Giới tính của thành viên & Nam, Nữ, Khác \\
        \rowcolor{tablerowalt}
        \texttt{NgaySinhDuongLich} & $\circ$ & Ngày sinh theo Dương lịch & 1995-10-20 \\
        \texttt{NgaySinhAmLich} & $\circ$ & Ngày sinh theo Âm lịch & Ngày 27 tháng 8 năm Ất Hợi \\
        \rowcolor{tablerowalt}
        \texttt{TheHe} & $*$ & Đời thứ mấy trong gia phả tính từ cụ Thủy tổ & 12 \\
        \texttt{ConThu} & $\circ$ & Thứ bậc sinh trong gia đình & Con trưởng, Con thứ ba \\
        \rowcolor{tablerowalt}
        \texttt{TinhTrang} & $*$ & Tình trạng sinh học & Còn sống, Đã mất \\
        \texttt{NgayMat} & $\circ$ & Ngày mất để nhắc ngày giỗ & 2020-04-15 (23/3 Âm lịch) \\
        \rowcolor{tablerowalt}
        \texttt{NoiAnTang} & $\circ$ & Vị trí mộ phần / nơi an táng & Nghĩa trang gia tộc Đô Lương \\
        \texttt{DiaChiHienTai} & $\circ$ & Nơi ở / địa chỉ cư trú hiện tại & Quận 1, TP. Hồ Chí Minh \\
        \rowcolor{tablerowalt}
        \texttt{AnhDaiDien} & $\circ$ & Đường dẫn ảnh chân dung của thành viên & \texttt{avatar\_001.jpg} \\
        \texttt{Tuoi} & Derived & Tuổi tính từ ngày sinh đến hiện tại & 31 \\
        \bottomrule
    \end{tabularx}
    \caption{Danh mục thuộc tính thực thể THANH\_VIEN (Member / Profile)}
    \label{tab:ent_thanh_vien}
\end{table}

\subsubsection{Thực thể HO\_SO\_NGHE\_NGHIEP (Career \& Education)}
Ghi nhận thông tin học vấn, chuyên môn và công việc của các thành viên phục vụ cho việc tạo Danh bạ gia đình (Family Directory), tìm kiếm và kết nối tương trợ giữa các thế hệ.

\begin{table}[H]
    \centering
    \small
    \begin{tabularx}{\textwidth}{l c X p{4.2cm}}
        \toprule
        \rowcolor{tableheader}\textcolor{white}{\textbf{Tên thuộc tính}} & \textcolor{white}{\textbf{Ký hiệu}} & \textcolor{white}{\textbf{Mô tả ý nghĩa nghiệp vụ}} & \textcolor{white}{\textbf{Ví dụ minh họa}} \\
        \midrule
        \rowcolor{tablerowalt}
        \texttt{MaHoSo} & $*$ & Mã hồ sơ chuyên môn danh bạ & \texttt{PROF\_001} \\
        \texttt{TrinhDoHocVan} & $\circ$ & Học vị / Bằng cấp cao nhất & Thạc sĩ, Cử nhân, Tiến sĩ \\
        \rowcolor{tablerowalt}
        \texttt{ChuyenNganh} & $\circ$ & Lĩnh vực chuyên môn đào tạo & Công nghệ thông tin, Y khoa \\
        \texttt{NgheNghiep} & $\circ$ & Chức danh nghề nghiệp hiện tại & Kỹ sư phần mềm, Bác sĩ \\
        \rowcolor{tablerowalt}
        \texttt{CoQuanCongTac} & $\circ$ & Nơi làm việc / Tổ chức công tác & FPT Software, BV Chợ Rẫy \\
        \texttt{LinhVucKinhDoanh} & $\circ$ & Lĩnh vực hoạt động doanh nghiệp & Bất động sản, Nông sản sạch \\
        \rowcolor{tablerowalt}
        \texttt{KhaNangHoTro} & $\circ$ & Lĩnh vực có thể chia sẻ, giúp đỡ con cháu & Tư vấn du học, Hướng nghiệp CNTT \\
        \bottomrule
    \end{tabularx}
    \caption{Danh mục thuộc tính thực thể HO\_SO\_NGHE\_NGHIEP (Career \& Education)}
    \label{tab:ent_ho_so_nghe_nghiep}
\end{table}

% ==============================================================================
% 3. NHÓM TƯƠNG TÁC & SỰ KIỆN
% ==============================================================================
\subsection{Nhóm Thực thể Tương tác \& Sự kiện (Community \& Events)}
Nhóm thực thể phục vụ mạng xã hội nội bộ dòng họ, cho phép chia sẻ tin tức, kết nối giao lưu và tổ chức các sự kiện chung.

\subsubsection{Thực thể BAI\_VIET (Post)}
Bài viết dạng tin tức dòng họ, thông báo sinh hoạt, chúc mừng hoặc câu chuyện gia đình được chia sẻ trên bảng tin cộng đồng.

\begin{table}[H]
    \centering
    \small
    \begin{tabularx}{\textwidth}{l c X p{4.2cm}}
        \toprule
        \rowcolor{tableheader}\textcolor{white}{\textbf{Tên thuộc tính}} & \textcolor{white}{\textbf{Ký hiệu}} & \textcolor{white}{\textbf{Mô tả ý nghĩa nghiệp vụ}} & \textcolor{white}{\textbf{Ví dụ minh họa}} \\
        \midrule
        \rowcolor{tablerowalt}
        \texttt{MaBaiViet} & $*$ & Mã định danh bài viết & \texttt{POST\_2026\_001} \\
        \texttt{TieuDe} & $*$ & Tiêu đề của bài viết & "Thông báo lễ mừng thọ đầu xuân" \\
        \rowcolor{tablerowalt}
        \texttt{NoiDung} & $*$ & Nội dung chi tiết bài viết & "Ban liên lạc trân trọng kính mời..." \\
        \texttt{PhamViChiaSe} & $*$ & Chế độ hiển thị & Toàn dòng họ, Toàn chi tộc, Gia đình nhỏ \\
        \rowcolor{tablerowalt}
        \texttt{LoaiBaiViet} & $*$ & Phân loại bài viết & Thông báo chính thức, Tin mừng/Hỷ, Tang chế \\
        \texttt{ThoiGianDang} & $*$ & Thời điểm đăng tải bài viết & 2026-02-10 08:00:00 \\
        \rowcolor{tablerowalt}
        \texttt{TrangThaiDuyet} & $*$ & Trạng thái kiểm duyệt & Đã duyệt, Chờ duyệt \\
        \bottomrule
    \end{tabularx}
    \caption{Danh mục thuộc tính thực thể BAI\_VIET (Post)}
    \label{tab:ent_bai_viet}
\end{table}

\subsubsection{Thực thể BINH\_LUAN (Comment)}
Ý kiến phản hồi, trao đổi của thành viên bên dưới mỗi bài viết hoặc sự kiện.

\begin{table}[H]
    \centering
    \small
    \begin{tabularx}{\textwidth}{l c X p{4.2cm}}
        \toprule
        \rowcolor{tableheader}\textcolor{white}{\textbf{Tên thuộc tính}} & \textcolor{white}{\textbf{Ký hiệu}} & \textcolor{white}{\textbf{Mô tả ý nghĩa nghiệp vụ}} & \textcolor{white}{\textbf{Ví dụ minh họa}} \\
        \midrule
        \rowcolor{tablerowalt}
        \texttt{MaBinhLuan} & $*$ & Mã định danh bình luận & \texttt{CMT\_00452} \\
        \texttt{NoiDung} & $*$ & Nội dung bình luận trao đổi & "Chi Hai xin đăng ký 5 suất tham dự ạ!" \\
        \rowcolor{tablerowalt}
        \texttt{ThoiGianTao} & $*$ & Thời điểm gửi bình luận & 2026-02-10 09:30:00 \\
        \texttt{TrangThai} & $*$ & Trạng thái hiển thị & Hiển thị, Đã ẩn \\
        \bottomrule
    \end{tabularx}
    \caption{Danh mục thuộc tính thực thể BINH\_LUAN (Comment)}
    \label{tab:ent_binh_luan}
\end{table}

\subsubsection{Thực thể ALBUM\_ANH (Album)}
Kho lưu trữ bộ sưu tập hình ảnh, video tư liệu về các hoạt động, buổi sum họp hoặc tài liệu lịch sử gia đình.

\begin{table}[H]
    \centering
    \small
    \begin{tabularx}{\textwidth}{l c X p{4.2cm}}
        \toprule
        \rowcolor{tableheader}\textcolor{white}{\textbf{Tên thuộc tính}} & \textcolor{white}{\textbf{Ký hiệu}} & \textcolor{white}{\textbf{Mô tả ý nghĩa nghiệp vụ}} & \textcolor{white}{\textbf{Ví dụ minh họa}} \\
        \midrule
        \rowcolor{tablerowalt}
        \texttt{MaAlbum} & $*$ & Mã định danh album & \texttt{ALB\_2026\_TET} \\
        \texttt{TenAlbum} & $*$ & Tên bộ ảnh / video & "Lễ Giỗ tổ Bính Ngọ 2026" \\
        \rowcolor{tablerowalt}
        \texttt{MoTa} & $\circ$ & Tóm tắt bối cảnh và ý nghĩa & Tập hợp ảnh sum họp từ đường đầu xuân \\
        \texttt{NgayTao} & $*$ & Thời điểm tạo album & 2026-02-15 14:00:00 \\
        \rowcolor{tablerowalt}
        \texttt{SoLuongAnh} & Derived & Tổng số bức ảnh có trong album & 45 \\
        \bottomrule
    \end{tabularx}
    \caption{Danh mục thuộc tính thực thể ALBUM\_ANH (Album)}
    \label{tab:ent_album_anh}
\end{table}

\subsubsection{Thực thể SU\_KIEN (Event)}
Lưu thông tin về các cuộc gặp gỡ, lễ giỗ họ, đại hội gia tộc, mừng thọ hoặc chương trình khuyến học.

\begin{table}[H]
    \centering
    \small
    \begin{tabularx}{\textwidth}{l c X p{4.2cm}}
        \toprule
        \rowcolor{tableheader}\textcolor{white}{\textbf{Tên thuộc tính}} & \textcolor{white}{\textbf{Ký hiệu}} & \textcolor{white}{\textbf{Mô tả ý nghĩa nghiệp vụ}} & \textcolor{white}{\textbf{Ví dụ minh họa}} \\
        \midrule
        \rowcolor{tablerowalt}
        \texttt{MaSuKien} & $*$ & Mã định danh sự kiện & \texttt{EVT\_GIO\_TO\_2026} \\
        \texttt{TenSuKien} & $*$ & Tên sự kiện & Lễ Giỗ tổ Thủy tổ năm 2026 \\
        \rowcolor{tablerowalt}
        \texttt{MoTa} & $\circ$ & Chi tiết lịch trình và nội dung sự kiện & 08h00: Dâng hương; 11h00: Thụ lộc \\
        \texttt{ThoiGianBatDau} & $*$ & Thời điểm bắt đầu diễn ra & 2026-03-03 07:30:00 \\
        \rowcolor{tablerowalt}
        \texttt{ThoiGianKetThuc} & $\circ$ & Thời điểm kết thúc dự kiến & 2026-03-03 13:00:00 \\
        \texttt{DiaDiem} & $*$ & Địa điểm tổ chức (Trực tiếp / Trực tuyến) & Nhà thờ tổ dòng họ, Đô Lương \\
        \rowcolor{tablerowalt}
        \texttt{NganSachDuKien} & $\circ$ & Kinh phí dự toán tổ chức sự kiện & 25.000.000 VNĐ \\
        \texttt{TrangThai} & $*$ & Trạng thái sự kiện & Sắp diễn ra, Đã hoàn thành \\
        \bottomrule
    \end{tabularx}
    \caption{Danh mục thuộc tính thực thể SU\_KIEN (Event)}
    \label{tab:ent_su_kien}
\end{table}

\subsubsection{Thực thể DANG\_KY\_SU\_KIEN (RSVP)}
Quản lý việc xác nhận tham gia sự kiện của các thành viên và số lượng người đi cùng, hỗ trợ khâu chuẩn bị hậu cần.

\begin{table}[H]
    \centering
    \small
    \begin{tabularx}{\textwidth}{l c X p{4.2cm}}
        \toprule
        \rowcolor{tableheader}\textcolor{white}{\textbf{Tên thuộc tính}} & \textcolor{white}{\textbf{Ký hiệu}} & \textcolor{white}{\textbf{Mô tả ý nghĩa nghiệp vụ}} & \textcolor{white}{\textbf{Ví dụ minh họa}} \\
        \midrule
        \rowcolor{tablerowalt}
        \texttt{MaDangKy} & $*$ & Mã xác nhận đăng ký & \texttt{RSVP\_001} \\
        \texttt{TrangThaiThamGia} & $*$ & Trạng thái phản hồi của thành viên & Sẽ tham gia, Không thể tham gia \\
        \rowcolor{tablerowalt}
        \texttt{SoNguoiDiCung} & $*$ & Số lượng người đi kèm gia đình & 3 \\
        \texttt{GhiChu} & $\circ$ & Lời nhắn gửi ban tổ chức & Cần hỗ trợ xe đưa đón người cao tuổi \\
        \rowcolor{tablerowalt}
        \texttt{ThoiDiemPhanHoi} & $*$ & Thời gian gửi xác nhận & 2026-02-20 18:45:00 \\
        \bottomrule
    \end{tabularx}
    \caption{Danh mục thuộc tính thực thể DANG\_KY\_SU\_KIEN (RSVP)}
    \label{tab:ent_dang_ky_su_kien}
\end{table}

% ==============================================================================
% 4. NHÓM DI SẢN SỐ
% ==============================================================================
\subsection{Nhóm Thực thể Di sản số (Family Heritage)}
Nhóm thực thể chịu trách nhiệm số hóa, bảo tồn các giá trị phi vật thể, văn bản truyền thống và tri thức dòng họ, đồng thời cung cấp dữ liệu nền tảng cho Trợ lý Trí tuệ nhân tạo (AI Assistant) tra cứu và tóm tắt.

\subsubsection{Thực thể DI\_SAN\_LICH\_SU (Heritage / Document)}
Lưu trữ các văn bản, hiện vật số hóa như: sắc phong, gia phả cổ, hương ước, văn tự chữ Hán -- Nôm, câu đối, di chúc hoặc các giai thoại truyền khẩu.

\begin{table}[H]
    \centering
    \small
    \begin{tabularx}{\textwidth}{l c X p{4.2cm}}
        \toprule
        \rowcolor{tableheader}\textcolor{white}{\textbf{Tên thuộc tính}} & \textcolor{white}{\textbf{Ký hiệu}} & \textcolor{white}{\textbf{Mô tả ý nghĩa nghiệp vụ}} & \textcolor{white}{\textbf{Ví dụ minh họa}} \\
        \midrule
        \rowcolor{tablerowalt}
        \texttt{MaDiSan} & $*$ & Mã định danh di sản/tư liệu & \texttt{HERITAGE\_01} \\
        \texttt{TieuDe} & $*$ & Tên gọi của tư liệu/di sản & "Sắc phong thời Cảnh Hưng năm 1783" \\
        \rowcolor{tablerowalt}
        \texttt{LoaiDiSan} & $*$ & Thể loại di sản số & Văn bản chữ Hán/Nôm, Điển tích \\
        \texttt{NienDai} & $\circ$ & Thời kỳ hoặc năm xuất hiện & Năm 1783 (Thời Hậu Lê) \\
        \rowcolor{tablerowalt}
        \texttt{NoiDungDichNghia} & $\circ$ & Nội dung phiên âm, dịch nghĩa Quốc ngữ & "Sắc phong cho cụ Thủy tổ..." \\
        \texttt{DuongDanTepTin} & $*$ & Đường dẫn tệp số hóa độ nét cao & \texttt{sac\_phong\_1783.pdf} \\
        \rowcolor{tablerowalt}
        \texttt{TomTatAI} & $\circ$ & Đoạn tóm tắt tạo bởi mô hình LLM & Ghi nhận công lao hộ quốc an dân... \\
        \texttt{GiaTriLichSu} & $\circ$ & Ý nghĩa giáo dục, giá trị văn hóa & Minh chứng truyền thống trung quân ái quốc \\
        \bottomrule
    \end{tabularx}
    \caption{Danh mục thuộc tính thực thể DI\_SAN\_LICH\_SU (Heritage / Document)}
    \label{tab:ent_di_san}
\end{table}

\subsubsection{Thực thể NHAN\_VAT\_TIEU\_BIEU (Notable Member)}
Lưu danh những nhân vật có công lao to lớn với dòng họ, xã hội hoặc các danh nhân văn hóa, anh hùng liệt sĩ, thủ khoa đỗ đạt để vinh danh và làm gương cho thế hệ trẻ.

\begin{table}[H]
    \centering
    \small
    \begin{tabularx}{\textwidth}{l c X p{4.2cm}}
        \toprule
        \rowcolor{tableheader}\textcolor{white}{\textbf{Tên thuộc tính}} & \textcolor{white}{\textbf{Ký hiệu}} & \textcolor{white}{\textbf{Mô tả ý nghĩa nghiệp vụ}} & \textcolor{white}{\textbf{Ví dụ minh họa}} \\
        \midrule
        \rowcolor{tablerowalt}
        \texttt{MaNhanVat} & $*$ & Mã hồ sơ vinh danh & \texttt{NOTABLE\_01} \\
        \texttt{DanhHieu} & $*$ & Danh xưng hoặc tước vị vinh danh & "Tiến sĩ khoa Canh Thìn 1460", "Nhà giáo ưu tú" \\
        \rowcolor{tablerowalt}
        \texttt{TieuSuChiTiet} & $*$ & Bối cảnh cuộc đời và sự nghiệp & Đỗ Đệ nhị giáp Tiến sĩ năm 1460... \\
        \texttt{CongLaoDongHo} & $\circ$ & Đóng góp cụ thể cho gia tộc & Hưng công xây dựng từ đường \\
        \rowcolor{tablerowalt}
        \texttt{CongHienXaHoi} & $\circ$ & Đóng góp đối với quê hương đất nước & Làm quan Thượng thư Bộ Lễ \\
        \texttt{TaiLieuThamKhao} & $\circ$ & Nguồn dẫn chứng, sách báo ghi chép & Đại Việt Sử ký Toàn thư, Văn bia \\
        \bottomrule
    \end{tabularx}
    \caption{Danh mục thuộc tính thực thể NHAN\_VAT\_TIEU\_BIEU (Notable Member)}
    \label{tab:ent_nhan_vat_tieu_bieu}
\end{table}

% !TEX root = ../../main.tex

\section[Xác định các mối quan hệ và bản số (Cardinality)]{Xác định các mối quan hệ (Relationships) và bản số (Cardinality)}
\label{sec:1_2_quan_he_cardinality}

Trong quá trình thiết kế cơ sở dữ liệu cho hệ thống \textbf{FamilyConnect}, việc xác định các mối quan hệ giữa các thực thể đóng vai trò quan trọng nhằm mô tả chính xác cách thức các đối tượng nghiệp vụ liên kết và tương tác với nhau. Các mối quan hệ được xây dựng dựa trên yêu cầu quản lý thành viên gia đình, cấu trúc gia phả, bảo mật phân quyền, hoạt động cộng đồng và chia sẻ di sản số.

Mỗi mối quan hệ cần được xác định rõ tên gọi, các thực thể tham gia, bản số (\textit{Cardinality}) và ràng buộc tham gia (\textit{Participation Constraint}). Việc phân tích này là cơ sở trực tiếp để chuyển đổi mô hình thực thể kết hợp (ER/EER Model) sang mô hình quan hệ ở các bước thiết kế tiếp theo.

% ==============================================================================
% 1. QUAN HỆ PHẢ HỆ CỐT LÕI
% ==============================================================================
\subsection{Quan hệ Cha/Mẹ -- Con (Quan hệ phụ mẫu)}

Quan hệ Cha/Mẹ -- Con biểu diễn mối liên kết huyết thống trực tiếp giữa các thế hệ thành viên trong gia đình. Đây là một mối quan hệ đệ quy nhị phân (\textit{Binary Recursive Relationship}), trong đó cùng một thực thể \texttt{THANH\_VIEN} tham gia với hai vai trò khác nhau: Cha/Mẹ (\textit{Parent}) và Con (\textit{Child}).

Mối quan hệ này cho phép hệ thống xây dựng và quản lý đồ thị cây gia phả có hướng, thể hiện dòng chảy thế hệ xuyên suốt lịch sử dòng họ.

\paragraph{Đặc điểm của mối quan hệ:}
\begin{itemize}[noitemsep,topsep=2pt]
    \item \textbf{Loại quan hệ:} Quan hệ đệ quy nhị phân (Binary Recursive Relationship).
    \item \textbf{Bản số (Cardinality):} Một Cha/Mẹ có thể có nhiều Con ($1:N$). Mỗi người con có tối đa một Cha ruột và tối đa một Mẹ ruột trong cơ sở dữ liệu.
    \item \textbf{Ràng buộc tham gia (Participation Constraint):}
    \begin{itemize}[noitemsep]
        \item Thực thể Con tham gia một phần (\textit{Partial Participation}), do các cụ Tiền nhân/Thủy tổ chưa rõ thông tin phụ mẫu.
        \item Thực thể Cha/Mẹ tham gia một phần (\textit{Partial Participation}), vì không phải mọi thành viên đều đã có con cái.
    \end{itemize}
\end{itemize}

\paragraph{Biểu diễn trong mô hình quan hệ:}
Có thể triển khai quan hệ này thông qua hai khóa ngoại đệ quy trực tiếp trong lược đồ quan hệ \texttt{THANH\_VIEN}:
\[
\textbf{THANH\_VIEN}(
\underline{\texttt{MaThanhVien}},
\texttt{HoVaTen},
\texttt{GioiTinh},
\texttt{MaCha},
\texttt{MaMe},
\dots)
\]
Trong đó, \texttt{MaCha} và \texttt{MaMe} là các khóa ngoại tự tham chiếu đến khóa chính \texttt{MaThanhVien} của chính bảng \texttt{THANH\_VIEN}.

% ==============================================================================
% 2. QUAN HỆ HÔN PHỐI
% ==============================================================================
\subsection{Quan hệ Hôn phối / Vợ chồng}

Quan hệ Hôn phối thể hiện mối liên kết hợp pháp hoặc tập quán giữa hai thành viên trong cộng đồng gia tộc. Khác với quan hệ phụ mẫu cố định suốt đời, quan hệ hôn phối có tính lịch sử, trạng thái thay đổi theo thời gian và mang theo các thuộc tính riêng (ngày cưới, ngày ly hôn, tình trạng).

\paragraph{Đặc điểm của mối quan hệ:}
\begin{itemize}[noitemsep,topsep=2pt]
    \item \textbf{Loại quan hệ:} Quan hệ đệ quy giữa hai cá nhân thuộc thực thể \texttt{THANH\_VIEN}.
    \item \textbf{Bản số (Cardinality):} Tại một thời điểm xác định, quan hệ hôn nhân là $1:1$. Tuy nhiên, xét trên toàn bộ dòng lịch sử, một cá nhân có thể trải qua nhiều cuộc hôn phối ($N:M$).
    \item \textbf{Thuộc tính của mối quan hệ:} \texttt{NgayKetHon}, \texttt{NgayLyHon}, \texttt{TrangThaiHonNhan}.
    \item \textbf{Ràng buộc tham gia:} Cả hai phía đều tham gia một phần, vì các thành viên chưa lập gia đình sẽ không xuất hiện trong quan hệ này.
\end{itemize}

\paragraph{Biểu diễn trong mô hình quan hệ:}
Do mối quan hệ có thuộc tính riêng và bản số lịch sử là $N:M$, quan hệ này được tách thành thực thể liên kết \texttt{HON\_NHAN}:
\[
\textbf{HON\_NHAN}(
\underline{\texttt{MaHonNhan}},
\texttt{MaChong},
\texttt{MaVo},
\texttt{NgayKetHon},
\texttt{NgayLyHon},
\texttt{TrangThai})
\]
Các trường \texttt{MaChong}, \texttt{MaVo} là khóa ngoại tham chiếu đến \texttt{MaThanhVien} trong \texttt{THANH\_VIEN}.

% ==============================================================================
% 3. QUAN HỆ DÒNG HỌ - CHI TỘC - THÀNH VIÊN
% ==============================================================================
\subsection{Quan hệ Dòng họ -- Chi tộc -- Thành viên}

Mối quan hệ này biểu diễn cơ cấu tổ chức phân cấp hành chính trong dòng họ truyền thống Việt Nam: một đại Dòng họ (\texttt{DONG\_HO}) gồm nhiều Chi tộc (\texttt{CHI\_TOC}), và mỗi Chi tộc quản lý nhiều Thành viên (\texttt{THANH\_VIEN}).

\paragraph{Đặc điểm của mối quan hệ:}
\begin{itemize}[noitemsep,topsep=2pt]
    \item \textbf{Quan hệ Dòng họ -- Chi tộc:}
    \begin{itemize}[noitemsep]
        \item Bản số: $1:N$. Một Dòng họ có nhiều Chi tộc; mỗi Chi tộc trực thuộc duy nhất một Dòng họ.
        \item Ràng buộc tham gia: Thực thể \texttt{CHI\_TOC} tham gia toàn bộ (\textit{Total}); \texttt{DONG\_HO} tham gia một phần (\textit{Partial}).
    \end{itemize}
    \item \textbf{Quan hệ Chi tộc -- Thành viên:}
    \begin{itemize}[noitemsep]
        \item Bản số: $1:N$. Một Chi tộc bao gồm nhiều Thành viên; mỗi Thành viên sinh ra hoặc nhập tịch vào một Chi tộc xác định.
        \item Ràng buộc tham gia: Thực thể \texttt{THANH\_VIEN} tham gia toàn bộ (\textit{Total}) để đảm bảo tính phân cấp; \texttt{CHI\_TOC} tham gia một phần.
    \end{itemize}
\end{itemize}

\paragraph{Biểu diễn trong mô hình quan hệ:}
Khóa ngoại \texttt{MaDongHo} được đặt trong bảng \texttt{CHI\_TOC}, và khóa ngoại \texttt{MaChiToc} được đặt trong bảng \texttt{THANH\_VIEN}:
\[
\textbf{CHI\_TOC}(
\underline{\texttt{MaChiToc}},
\texttt{TenChiToc},
\texttt{DoiThu},
\texttt{MaDongHo})
\]
\[
\textbf{THANH\_VIEN}(
\underline{\texttt{MaThanhVien}},
\texttt{HoVaTen},
\texttt{MaChiToc},
\dots)
\]

% ==============================================================================
% 4. QUAN HỆ NGƯỜI DÙNG, BẢO MẬT & HỒ SƠ
% ==============================================================================
\subsection{Quan hệ Người dùng, Phân quyền và Hồ sơ chuyên môn}

Hệ thống kết hợp giữa tài khoản trực tuyến (\texttt{NGUOI\_DUNG}) và hồ sơ phả hệ (\texttt{THANH\_VIEN}), đồng thời kiểm soát quyền hạn theo mô hình RBAC:

\begin{itemize}[noitemsep,topsep=2pt]
    \item \textbf{Liên kết Tài khoản -- Hồ sơ (\textit{Profile Claiming}):} Mối quan hệ $1:1$ tùy chọn giữa \texttt{NGUOI\_DUNG} và \texttt{THANH\_VIEN}. Một người dùng có thể liên kết tối đa với một hồ sơ thành viên trong gia phả và ngược lại.
    \item \textbf{Gán Vai trò \& Phân quyền (\texttt{NGUOI\_DUNG} -- \texttt{VAI\_TRO}):} Bản số $N:M$. Một tài khoản đảm nhiệm nhiều vai trò (Quản trị dòng họ, Trưởng chi tộc, Thành viên thường). Quan hệ được thực thể hóa qua bảng trung gian \texttt{NGUOI\_DUNG\_VAI\_TRO}.
    \item \textbf{Ma trận Quyền hạn (\texttt{VAI\_TRO} -- \texttt{QUYEN}):} Bản số $N:M$, xác định tập hành động nghiệp vụ được phép cho từng vai trò. Quan hệ được thực thể hóa qua bảng trung gian \texttt{VAI\_TRO\_QUYEN}:
\end{itemize}
\[
\textbf{VAI\_TRO\_QUYEN}(
\underline{\texttt{MaVaiTro}},
\underline{\texttt{MaQuyen}})
\]
\begin{itemize}[noitemsep,topsep=2pt]
    \item \textbf{Hồ sơ Nghề nghiệp (\texttt{THANH\_VIEN} -- \texttt{HO\_SO\_NGHE\_NGHIEP}):} Quan hệ $1:1$, mở rộng thông tin học vấn, chuyên môn và nghề nghiệp của thành viên phục vụ Danh bạ gia tộc.
\end{itemize}

% ==============================================================================
% 5. QUAN HỆ BÀI VIẾT & TƯƠNG TÁC CỘNG ĐỒNG
% ==============================================================================
\subsection{Quan hệ Bài viết và Tương tác Cộng đồng}

Phân hệ mạng xã hội nội bộ cho phép thành viên chia sẻ thông tin, lưu giữ khoảnh khắc và thảo luận:

\begin{itemize}[noitemsep,topsep=2pt]
    \item \textbf{Thành viên tạo Bài viết (\texttt{THANH\_VIEN} -- \texttt{BAI\_VIET}):} Bản số $1:N$. Mỗi bài viết bắt buộc phải có một tác giả khởi tạo (\texttt{MaNguoiTao}), do đó \texttt{BAI\_VIET} tham gia toàn bộ.
    \item \textbf{Bài viết và Bình luận (\texttt{BAI\_VIET} -- \texttt{BINH\_LUAN}):} Bản số $1:N$. Mỗi bình luận gắn với đúng một bài viết.
    \item \textbf{Thành viên gửi Bình luận (\texttt{THANH\_VIEN} -- \texttt{BINH\_LUAN}):} Bản số $1:N$. Thành viên gửi ý kiến trao đổi dưới các bài đăng.
    \item \textbf{Thành viên tạo Album ảnh (\texttt{THANH\_VIEN} -- \texttt{ALBUM\_ANH}):} Bản số $1:N$. Lưu trữ các bộ sưu tập hình ảnh sự kiện và tư liệu gia đình.
\end{itemize}

% ==============================================================================
% 6. QUAN HỆ SỰ KIỆN & ĐĂNG KÝ THAM GIA (RSVP)
% ==============================================================================
\subsection{Quan hệ Sự kiện và Đăng ký Tham gia}

Nhằm điều phối các buổi lễ giỗ tổ, họp họ, mừng thọ và đại hội thường niên:

\begin{itemize}[noitemsep,topsep=2pt]
    \item \textbf{Tổ chức Sự kiện (\texttt{THANH\_VIEN} -- \texttt{SU\_KIEN}):} Bản số $1:N$. Một thành viên (Ban tổ chức/Trưởng tộc) đứng ra tạo và quản lý sự kiện.
    \item \textbf{Đăng ký Tham gia (\texttt{THANH\_VIEN} -- \texttt{SU\_KIEN}):} Bản số $N:M$. Một sự kiện có nhiều thành viên tham gia, và một thành viên có thể tham dự nhiều sự kiện. Mối quan hệ được thực thể hóa qua thực thể liên kết \texttt{DANG\_KY\_SU\_KIEN} (RSVP) chứa các thông tin như số lượng người đi kèm, trạng thái phản hồi và ghi chú hậu cần:
\end{itemize}
\[
\begin{aligned}
\textbf{DANG\_KY\_SU\_KIEN}(
&\underline{\texttt{MaDangKy}},
\texttt{MaThanhVien},
\texttt{MaSuKien}, \\
&\texttt{TrangThaiThamGia},
\texttt{SoNguoiDiCung})
\end{aligned}
\]

% ==============================================================================
% 7. QUAN HỆ DI SẢN & TÔN VINH NHÂN VẬT
% ==============================================================================
\subsection{Quan hệ Di sản số và Tôn vinh Nhân vật tiêu biểu}

Phân hệ văn hóa di sản lưu trữ và số hóa các giá trị tinh thần của gia tộc:

\begin{itemize}[noitemsep,topsep=2pt]
    \item \textbf{Đóng góp Di sản (\texttt{THANH\_VIEN} -- \texttt{DI\_SAN\_LICH\_SU}):} Bản số $1:N$ (hoặc $N:M$ khi nhiều thành viên cùng đóng góp/số hóa). Mỗi tài liệu, sắc phong, văn tự cổ được bảo trợ hoặc đăng tải bởi thành viên dòng họ.
    \item \textbf{Tôn vinh Nhân vật tiêu biểu (\texttt{THANH\_VIEN} -- \texttt{NHAN\_VAT\_TIEU\_BIEU}):} Bản số $1:1$ (hoặc $1:N$ theo thời kỳ vinh danh). Mỗi hồ sơ danh nhân tôn vinh chính xác một tiền nhân hoặc cá nhân kiệt xuất trong cây gia phả.
\end{itemize}

% ==============================================================================
% 8. QUAN HỆ NGƯỜI DÙNG & NHẬT KÝ HỆ THỐNG
% ==============================================================================
\subsection{Quan hệ Người dùng và Nhật ký Hệ thống}

Mọi thao tác nhạy cảm trên dữ liệu (thêm/sửa/xóa phả hệ, duyệt bài viết, phân quyền) đều được ghi lại tự động bởi hệ thống:

\begin{itemize}[noitemsep,topsep=2pt]
    \item \textbf{Người dùng tạo Nhật ký (\texttt{NGUOI\_DUNG} -- \texttt{NHAT\_KY\_HE\_THONG}):} Bản số $1:N$. Một tài khoản có thể phát sinh nhiều bản ghi nhật ký; mỗi bản ghi nhật ký gắn liền với đúng một người dùng thực hiện thao tác. Thực thể \texttt{NHAT\_KY\_HE\_THONG} tham gia toàn bộ -- mọi bản ghi đều có chủ thể tác động.
    \item \textbf{Phạm vi theo dõi:} Nhật ký theo dõi tất cả các thực thể cốt lõi: \texttt{THANH\_VIEN}, \texttt{BAI\_VIET}, \texttt{SU\_KIEN}, \texttt{DI\_SAN\_LICH\_SU} và cấu hình phân quyền. Trường \texttt{DoiTuongTacDong} lưu tên thực thể bị tác động, \texttt{DuLieuCu}/\texttt{DuLieuMoi} lưu snapshot trạng thái trước/sau.
\end{itemize}

% ==============================================================================
% 8. MA TRẬN TỔNG HỢP BẢN SỐ (CARDINALITY MATRIX)
% ==============================================================================
\subsection{Bảng tổng hợp bản số (Cardinality Matrix)}

Bảng~\ref{tab:cardinality_matrix} tổng hợp toàn diện các mối quan hệ nghiệp vụ chính giữa 15 thực thể cốt lõi trong hệ thống \textbf{FamilyConnect}:

\begin{table}[H]
    \centering
    \small
    \renewcommand{\arraystretch}{1.25}
    \begin{tabularx}{\textwidth}{p{3.8cm} p{4.2cm} c X}
        \toprule
        \rowcolor{tableheader}\textcolor{white}{\textbf{Mối quan hệ}} & \textcolor{white}{\textbf{Các thực thể tham gia}} & \textcolor{white}{\textbf{Bản số}} & \textcolor{white}{\textbf{Mô tả \& Ràng buộc tham gia}} \\
        \midrule
        \rowcolor{tablerowalt}
        Phụ mẫu -- Con cái & \texttt{THANH\_VIEN} -- \texttt{THANH\_VIEN} & $1:N$ & Quan hệ đệ quy huyết thống (Cha/Mẹ ruột). Con tham gia một phần. \\
        Hôn phối / Vợ chồng & \texttt{THANH\_VIEN} -- \texttt{THANH\_VIEN} & $N:M$ & Thực thể hóa qua \texttt{HON\_NHAN}. Lưu lịch sử và tình trạng hôn nhân. \\
        \rowcolor{tablerowalt}
        Dòng họ -- Chi tộc & \texttt{DONG\_HO} -- \texttt{CHI\_TOC} & $1:N$ & Phân cấp dòng họ. \texttt{CHI\_TOC} tham gia toàn bộ. \\
        Chi tộc -- Thành viên & \texttt{CHI\_TOC} -- \texttt{THANH\_VIEN} & $1:N$ & Quản lý nhân khẩu. \texttt{THANH\_VIEN} tham gia toàn bộ. \\
        \rowcolor{tablerowalt}
        Tài khoản -- Hồ sơ & \texttt{NGUOI\_DUNG} -- \texttt{THANH\_VIEN} & $1:1$ & Liên kết tài khoản đăng nhập với hồ sơ gia phả (Tùy chọn). \\
        Phân quyền vai trò & \texttt{NGUOI\_DUNG} -- \texttt{VAI\_TRO} & $N:M$ & Thực thể hóa qua \texttt{NGUOI\_DUNG\_VAI\_TRO}. Gán vai trò cho người dùng. \\
        \rowcolor{tablerowalt}
        Vai trò -- Quyền hạn & \texttt{VAI\_TRO} -- \texttt{QUYEN} & $N:M$ & Thực thể hóa qua \texttt{VAI\_TRO\_QUYEN}. Ma trận phân quyền nghiệp vụ. \\
        Hồ sơ chuyên môn & \texttt{THANH\_VIEN} -- \texttt{HO\_SO\_NGHE\_NGHIEP} & $1:1$ & Mở rộng thông tin học vấn và nghề nghiệp thành viên. \\
        \rowcolor{tablerowalt}
        Tác giả bài viết & \texttt{THANH\_VIEN} -- \texttt{BAI\_VIET} & $1:N$ & Thành viên đăng tải tin tức, bài viết gia tộc. \\
        Bài viết -- Bình luận & \texttt{BAI\_VIET} -- \texttt{BINH\_LUAN} & $1:N$ & Một bài viết có nhiều bình luận tương tác. \\
        \rowcolor{tablerowalt}
        Tác giả bình luận & \texttt{THANH\_VIEN} -- \texttt{BINH\_LUAN} & $1:N$ & Thành viên gửi ý kiến thảo luận. \\
        Tạo Album tư liệu & \texttt{THANH\_VIEN} -- \texttt{ALBUM\_ANH} & $1:N$ & Thành viên lưu trữ bộ sưu tập hình ảnh, tư liệu gia đình. \\
        \rowcolor{tablerowalt}
        Chủ trì sự kiện & \texttt{THANH\_VIEN} -- \texttt{SU\_KIEN} & $1:N$ & Ban tổ chức tạo lịch sự kiện, lễ giỗ họ. \\
        Đăng ký sự kiện & \texttt{THANH\_VIEN} -- \texttt{SU\_KIEN} & $N:M$ & Thực thể hóa qua \texttt{DANG\_KY\_SU\_KIEN} (RSVP). \\
        \rowcolor{tablerowalt}
        Đóng góp di sản & \texttt{THANH\_VIEN} -- \texttt{DI\_SAN\_LICH\_SU} & $1:N$ & Thành viên đóng góp, số hóa tài liệu di sản. \\
        Tôn vinh danh nhân & \texttt{THANH\_VIEN} -- \texttt{NHAN\_VAT\_TIEU\_BIEU} & $1:1$ & Vinh danh nhân vật kiệt xuất trong dòng tộc. \\
        \rowcolor{tablerowalt}
        Ghi nhật ký & \texttt{NGUOI\_DUNG} -- \texttt{NHAT\_KY\_HE\_THONG} & $1:N$ & Audit trail tự động. Mọi bản ghi nhật ký có chủ thể tác động. \\
        \bottomrule
    \end{tabularx}
    \caption{Ma trận tổng hợp các mối quan hệ và bản số (Cardinality Matrix)}
    \label{tab:cardinality_matrix}
\end{table}

% ==============================================================================
% 9. NHẬN XÉT
% ==============================================================================
\subsection{Nhận xét và Đánh giá}

Qua quá trình phân tích và chuẩn hóa, mô hình quan niệm của nền tảng \textbf{FamilyConnect} phản ánh đầy đủ các đặc trưng cấu trúc sau:
\begin{enumerate}[noitemsep,topsep=2pt]
    \item \textbf{Sự đa dạng về bản số:} Hệ thống bao quát đầy đủ các dạng bản số: $1:N$, $1:1$ và $N:M$. Tổng cộng 17 mối quan hệ nghiệp vụ được xác định, trong đó 5 quan hệ $N:M$ (Hôn phối, Phân quyền vai trò, Ma trận quyền hạn, Đăng ký sự kiện) đều được thực thể hóa thành các bảng liên kết riêng.
    \item \textbf{Tính chất đệ quy đặc thù:} Quan hệ phụ mẫu (Cha/Mẹ -- Con) là trung tâm của cây phả hệ, đòi hỏi xử lý đệ quy chặt chẽ và tuân thủ các ràng buộc cấu trúc đồ thị không chu trình (DAG) theo RB03.
    \item \textbf{Tính toàn vẹn và phân cấp:} Ràng buộc tham gia toàn bộ của thành viên đối với chi tộc và dòng họ đảm bảo không tồn tại dữ liệu mồ côi (\textit{orphan records}), tạo nền tảng vững chắc cho việc thiết kế lược đồ vật lý và chuẩn hóa dữ liệu ở các chương tiếp theo.
    \item \textbf{Tính truy vết và kiểm toán:} Quan hệ Người dùng -- Nhật ký hệ thống đảm bảo mọi thao tác nghiệp vụ nhạy cảm đều có dấu vết truy xuất, hỗ trợ kiểm toán và khôi phục sự cố.
\end{enumerate}

% !TEX root = ../../main.tex
\section{Mô hình Thực thể -- Liên kết mở rộng (ERD / EERD)}
\label{sec:1_3_mo_hinh_erd_eerd}

% =====================================================
% GIỚI THIỆU
% =====================================================
\subsection{Giới thiệu}
% Mục này mô tả mục đích của tài liệu: trình bày mô hình thực thể (ERD)
% mở rộng (EERD) cho hệ thống FamilyConnect, bao gồm các thực thể chính,
% thuộc tính khóa, thuộc tính số (kiểu dữ liệu) và các mối quan hệ giữa chúng.
Tài liệu này trình bày mô hình thực thể -- mối quan hệ (ERD) mở rộng (EERD)
cho nền tảng \textbf{FamilyConnect}. Mô hình bao gồm 11 thực thể chính,
thể hiện đầy đủ thuộc tính khóa (PK), thuộc tính khóa ngoại (FK), thuộc tính
số liệu, và các mối quan hệ (kèm bậc số -- cardinality) giữa các thực thể.

% =====================================================
% DANH SÁCH THỰC THỂ
% =====================================================
\subsection{Danh sách thực thể và thuộc tính}

% ---- USER ----
\subsubsection{USER (Người dùng)}
% Thực thể đại diện cho tài khoản đăng nhập hệ thống, phục vụ nhóm
% chức năng "Người dùng & Bảo mật" (đăng ký, xác thực, RBAC).
\begin{longtable}{@{}p{3cm}p{2cm}p{9.5cm}@{}}
\toprule
\textbf{Thuộc tính} & \textbf{Khóa} & \textbf{Ghi chú} \\
\midrule
id & PK & Định danh duy nhất (uuid) \\
username & & Tên đăng nhập \\
email & & Địa chỉ email, dùng xác thực \\
password\_hash & & Mật khẩu đã mã hóa \\
role & & Vai trò (RBAC): quản trị viên / thành viên \\
\bottomrule
\end{longtable}

% ---- FAMILY ----
\subsubsection{FAMILY (Gia đình)}
% Thực thể gốc đại diện cho một gia đình / chi / nhánh, là điểm neo
% cho toàn bộ dữ liệu gia phả, cộng đồng, sự kiện và di sản.
\begin{longtable}{@{}p{3cm}p{2cm}p{9.5cm}@{}}
\toprule
\textbf{Thuộc tính} & \textbf{Khóa} & \textbf{Ghi chú} \\
\midrule
id & PK & Định danh duy nhất \\
name & & Tên gia đình / dòng họ \\
branch & & Chi / nhánh / phái \\
description & & Mô tả tổng quan \\
\bottomrule
\end{longtable}

% ---- PERSON ----
\subsubsection{PERSON (Thành viên gia phả)}
% Thực thể lõi của phân hệ gia phả. Thuộc tính parent_id tự tham chiếu
% (self-reference) đến chính thực thể PERSON để biểu diễn quan hệ
% cha/mẹ - con qua nhiều thế hệ mà không cần bảng riêng cho từng cấp.
\begin{longtable}{@{}p{3cm}p{2cm}p{9.5cm}@{}}
\toprule
\textbf{Thuộc tính} & \textbf{Khóa} & \textbf{Ghi chú} \\
\midrule
id & PK & Định danh duy nhất \\
family\_id & FK & Tham chiếu đến FAMILY \\
user\_id & FK & Tham chiếu đến USER (nếu thành viên có tài khoản) \\
parent\_id & FK & Tự tham chiếu đến PERSON (quan hệ cha/mẹ) \\
full\_name & & Họ và tên đầy đủ \\
dob & & Ngày sinh \\
dod & & Ngày mất (nếu có) \\
gender & & Giới tính \\
occupation & & Nghề nghiệp / học vấn \\
\bottomrule
\end{longtable}

% ---- MARRIAGE ----
\subsubsection{MARRIAGE (Hôn nhân)}
% Tách thành bảng riêng thay vì gắn trực tiếp vào PERSON, vì một người
% có thể có nhiều cuộc hôn nhân theo thời gian (đảm bảo chuẩn hóa dữ liệu).
\begin{longtable}{@{}p{3cm}p{2cm}p{9.5cm}@{}}
\toprule
\textbf{Thuộc tính} & \textbf{Khóa} & \textbf{Ghi chú} \\
\midrule
id & PK & Định danh duy nhất \\
person1\_id & FK & Tham chiếu đến PERSON (vợ/chồng 1) \\
person2\_id & FK & Tham chiếu đến PERSON (vợ/chồng 2) \\
marriage\_date & & Ngày kết hôn \\
\bottomrule
\end{longtable}

% ---- POST / COMMENT ----
\subsubsection{POST (Bài đăng) \& COMMENT (Bình luận)}
% Phục vụ nhóm chức năng "Cộng đồng & Tương tác": chia sẻ tin tức,
% hình ảnh gia đình và tương tác bình luận.
\begin{longtable}{@{}p{3cm}p{2cm}p{9.5cm}@{}}
\toprule
\textbf{Thuộc tính} & \textbf{Khóa} & \textbf{Ghi chú} \\
\midrule
\multicolumn{3}{l}{\textit{POST}} \\
id & PK & Định danh duy nhất \\
user\_id & FK & Người đăng bài \\
family\_id & FK & Bài đăng thuộc gia đình nào \\
content & & Nội dung bài đăng \\
created\_at & & Thời điểm đăng \\
\addlinespace
\multicolumn{3}{l}{\textit{COMMENT}} \\
id & PK & Định danh duy nhất \\
post\_id & FK & Tham chiếu đến POST \\
user\_id & FK & Người bình luận \\
content & & Nội dung bình luận \\
created\_at & & Thời điểm bình luận \\
\bottomrule
\end{longtable}

% ---- EVENT / RSVP ----
\subsubsection{EVENT (Sự kiện) \& RSVP}
% Phục vụ nhóm "Quản lý Sự kiện": tạo sự kiện, nhắc nhở, quản lý người
% tham gia. RSVP dùng khóa chính kép (event_id, user_id) -- không cần id riêng.
\begin{longtable}{@{}p{3cm}p{2cm}p{9.5cm}@{}}
\toprule
\textbf{Thuộc tính} & \textbf{Khóa} & \textbf{Ghi chú} \\
\midrule
\multicolumn{3}{l}{\textit{EVENT}} \\
id & PK & Định danh duy nhất \\
family\_id & FK & Sự kiện thuộc gia đình nào \\
name & & Tên sự kiện \\
event\_date & & Ngày diễn ra \\
location & & Địa điểm \\
\addlinespace
\multicolumn{3}{l}{\textit{RSVP}} \\
event\_id & PK, FK & Khóa chính kép -- Tham chiếu EVENT \\
user\_id & PK, FK & Khóa chính kép -- Tham chiếu USER \\
status & & Trạng thái tham gia (đồng ý / từ chối / chưa rõ) \\
\bottomrule
\end{longtable}

% ---- DOCUMENT ----
\subsubsection{DOCUMENT (Tài liệu / Di sản)}
% Phục vụ nhóm "Danh bạ & Di sản": lưu trữ tài liệu, câu chuyện lịch sử,
% lưu trữ số của gia đình.
\begin{longtable}{@{}p{3cm}p{2cm}p{9.5cm}@{}}
\toprule
\textbf{Thuộc tính} & \textbf{Khóa} & \textbf{Ghi chú} \\
\midrule
id & PK & Định danh duy nhất \\
family\_id & FK & Tài liệu thuộc gia đình nào \\
title & & Tiêu đề tài liệu \\
file\_url & & Đường dẫn lưu trữ tệp \\
story\_text & & Nội dung câu chuyện lịch sử \\
\bottomrule
\end{longtable}

% ---- NOTIFICATION ----
\subsubsection{NOTIFICATION (Thông báo)}
% Phục vụ hệ thống thông báo trong nhóm "Cộng đồng & Tương tác".
\begin{longtable}{@{}p{3cm}p{2cm}p{9.5cm}@{}}
\toprule
\textbf{Thuộc tính} & \textbf{Khóa} & \textbf{Ghi chú} \\
\midrule
id & PK & Định danh duy nhất \\
user\_id & FK & Người nhận thông báo \\
content & & Nội dung thông báo \\
is\_read & & Trạng thái đã đọc / chưa đọc \\
\bottomrule
\end{longtable}

% =====================================================
% MỐI QUAN HỆ
% =====================================================
\subsection{Danh sách mối quan hệ (Relationships)}
% Bảng tổng hợp các mối quan hệ giữa các thực thể, kèm bậc số
% (cardinality) theo ký hiệu chuẩn ERD: 1 - n (một - nhiều).
\begin{longtable}{@{}p{3.2cm}p{1.2cm}p{3.2cm}p{6.5cm}@{}}
\toprule
\textbf{Thực thể A} & \textbf{Bậc số} & \textbf{Thực thể B} & \textbf{Diễn giải} \\
\midrule
USER & 1--n & PERSON & Một tài khoản có thể liên kết một hồ sơ gia phả \\
FAMILY & 1--n & PERSON & Một gia đình gồm nhiều thành viên \\
PERSON & 1--n & PERSON & Quan hệ cha/mẹ -- con (tự tham chiếu) \\
PERSON & 1--n & MARRIAGE & Một người có thể có nhiều cuộc hôn nhân \\
USER & 1--n & POST & Một người dùng đăng nhiều bài viết \\
FAMILY & 1--n & POST & Bài viết thuộc về một gia đình \\
POST & 1--n & COMMENT & Một bài viết nhận nhiều bình luận \\
USER & 1--n & COMMENT & Một người dùng viết nhiều bình luận \\
FAMILY & 1--n & EVENT & Một gia đình tổ chức nhiều sự kiện \\
EVENT & 1--n & RSVP & Một sự kiện nhận nhiều phản hồi tham gia \\
USER & 1--n & RSVP & Một người dùng tham gia nhiều sự kiện \\
FAMILY & 1--n & DOCUMENT & Một gia đình lưu trữ nhiều tài liệu \\
USER & 1--n & NOTIFICATION & Một người dùng nhận nhiều thông báo \\
\bottomrule
\end{longtable}

% !TEX root = ../../main.tex
\section{Ràng buộc toàn vẹn mức quan niệm (Business Rules)}
\label{sec:1_4_rang_buoc_quan_niem}

Ràng buộc toàn vẹn mức quan niệm (Conceptual Integrity Constraints) là tập hợp các quy tắc ngữ nghĩa biểu diễn tính đúng đắn, hợp lệ và nhất quán của dữ liệu trong thế giới thực mà cơ sở dữ liệu phải thỏa mãn tại mọi thời điểm. Các ràng buộc này hoàn toàn độc lập với hệ quản trị cơ sở dữ liệu (DBMS) và phản ánh trung thực các quy luật sinh học tự nhiên, tập quán văn hóa dòng họ Việt Nam cùng các chính sách phân quyền đa cấp của nền tảng \textbf{FamilyConnect}.

Trong thiết kế hệ thống, mỗi ràng buộc toàn vẹn được chuẩn hóa gồm 4 thành phần:
\begin{itemize}[noitemsep,topsep=2pt]
    \item \textbf{Nội dung nghiệp vụ (Semantics):} Diễn giải quy tắc bằng ngôn ngữ tự nhiên.
    \item \textbf{Biểu diễn hình thức (Formal Specification):} Mô hình hóa bằng vị từ logic toán học.
    \item \textbf{Bối cảnh áp dụng (Context):} Xác định các thực thể và thuộc tính chịu ảnh hưởng.
    \item \textbf{Quy tắc toàn vẹn (Integrity Rule):} Điều kiện dữ liệu phải thỏa mãn tại mọi thời điểm.
\end{itemize}

% ==============================================================================
% 1.4.1. NHÓM QUY TẮC HUYẾT THỐNG VÀ PHẢ HỆ
% ==============================================================================
\subsection{Nhóm quy tắc toàn vẹn huyết thống và cấu trúc phả hệ}
Cây phả hệ là một đồ thị quan hệ huyết thống có hướng. Các thực thể \texttt{THANH\_VIEN} và các mối liên kết cha--mẹ, hôn phối bắt buộc phải tuân thủ các quy luật sinh học và cấu trúc đồ thị toán học sau:

% --- RB01 ---
\subsubsection{RB01 --- Tính duy nhất của quan hệ phụ mẫu ruột (Single Biological Parentage)}
\begin{itemize}
    \item \textbf{Nội dung nghiệp vụ:} Mỗi cá nhân trong gia phả chỉ có tối đa một người cha ruột và tối đa một người mẹ ruột. Cho phép giá trị rỗng (NULL) đối với tiền nhân chưa rõ căn cước, nhưng tuyệt đối không tồn tại hai cha ruột hoặc hai mẹ ruột.
    \item \textbf{Biểu diễn toán học:} Gọi $TV = \text{THANH\_VIEN}$. Đặt $\text{ChaRuot}(f, p)$, $\text{MeRuot}(m, p)$ là các vị từ quan hệ phụ mẫu:
    \begin{alignat}{2}
        &\forall p \in TV, \quad
        &&\big|\{ f \in TV \mid \text{ChaRuot}(f, p) \}\big| \le 1 \\[3pt]
        &\forall p \in TV, \quad
        &&\big|\{ m \in TV \mid \text{MeRuot}(m, p) \}\big| \le 1
    \end{alignat}
    \item \textbf{Bối cảnh:} Thực thể \texttt{THANH\_VIEN}, thuộc tính \texttt{MaCha}, \texttt{MaMe}, \texttt{GioiTinh}.
    \item \textbf{Quy tắc toàn vẹn:} Không được phép ghi nhận thêm quan hệ phụ mẫu nếu cá nhân đã có cha hoặc mẹ ruột cùng giới tính tương ứng.
\end{itemize}

% --- RB02 ---
\subsubsection{RB02 --- Ràng buộc mốc thời gian sinh học (Biological Age Consistency)}
\begin{itemize}
    \item \textbf{Nội dung nghiệp vụ:}
    \begin{enumerate}[label=(\alph*),noitemsep]
        \item Con cái phải sinh sau cha/mẹ ít nhất 15 năm (độ tuổi sinh sản tối thiểu).
        \item Người mẹ không thể sinh con sau ngày mất của chính mình.
        \item Người cha không thể có con sinh sau ngày mất của mình quá 300 ngày (giới hạn y học).
    \end{enumerate}
    \item \textbf{Biểu diễn toán học:} Viết tắt $\text{ns}(x) = x.\text{NgaySinh}$, $\text{nm}(x) = x.\text{NgayMat}$:
    \begin{alignat}{2}
        &\forall c, p \in TV, \quad
        &&\text{ChaMe}(p, c) \implies \text{ns}(c) \ge \text{ns}(p) + 15\text{ n\u0103m} \\[3pt]
        &\forall c, m \in TV, \quad
        &&\text{MeRuot}(m, c) \land \text{nm}(m) \ne \emptyset \implies \text{ns}(c) \le \text{nm}(m) \\[3pt]
        &\forall c, f \in TV, \quad
        &&\text{ChaRuot}(f, c) \land \text{nm}(f) \ne \emptyset \implies \text{ns}(c) \le \text{nm}(f) + 300\text{ ng\u00e0y}
    \end{alignat}
    \item \textbf{Bối cảnh:} Thực thể \texttt{THANH\_VIEN}, thuộc tính \texttt{NgaySinh\-DuongLich}, \texttt{NgayMat}, \texttt{GioiTinh}.
    \item \textbf{Quy tắc toàn vẹn:} Mọi bộ dữ liệu $(c, p)$ có quan hệ phụ mẫu đều phải thỏa mãn đồng thời cả ba điều kiện thời gian trên.
\end{itemize}

% --- RB03 ---
\subsubsection{RB03 --- Ràng buộc Đồ thị có hướng không chu trình (DAG Acyclicity)}
\begin{itemize}
    \item \textbf{Nội dung nghiệp vụ:} Cây phả hệ thể hiện dòng chảy thế hệ đơn hướng. Một cá nhân không thể vừa là tổ tiên vừa là hậu duệ của chính mình.
    \item \textbf{Biểu diễn toán học:} Gọi $G = (V, E)$ là đồ thị phả hệ với $V = TV$, $E = \{(u,v) \mid \text{ChaMe}(u,v)\}$. Gọi $E^{+}$ là bao đóng bắc cầu của $E$:
    \begin{equation}
        \forall p \in TV: \quad
        (p, p) \notin E^{+}
        \quad\iff\quad
        \text{Ancestors}(p) \cap \text{Descendants}(p) = \emptyset
    \end{equation}
    \item \textbf{Bối cảnh:} Quan hệ phụ mẫu giữa các thực thể \texttt{THANH\_VIEN}.
    \item \textbf{Quy tắc toàn vẹn:} Mọi thao tác thêm/sửa liên kết phụ mẫu đều phải kiểm tra chu trình trong đồ thị $G$ và bị từ chối nếu tạo ra vòng lặp.
\end{itemize}

% --- RB04 ---
\subsubsection{RB04 --- Tính tăng đơn điệu của bậc thế hệ (Generation Monotonicity)}
\begin{itemize}
    \item \textbf{Nội dung nghiệp vụ:} Bậc thế hệ của con cái luôn bằng bậc thế hệ của người cha cộng thêm 1 đơn vị (tính từ Cụ Khởi tổ là đời thứ 1).
    \item \textbf{Biểu diễn toán học:}
    \begin{equation}
        \forall c, f \in TV: \quad
        \text{ChaRuot}(f, c) \implies c.\text{TheHe} = f.\text{TheHe} + 1
    \end{equation}
    \item \textbf{Bối cảnh:} Thực thể \texttt{THANH\_VIEN}, thuộc tính \texttt{TheHe}, \texttt{MaCha}.
    \item \textbf{Quy tắc toàn vẹn:} Bậc thế hệ \texttt{TheHe} phải được tính toán nhất quán khi gắn nối thành viên vào cây phả hệ; không được phép gán bậc thế hệ tùy tiện.
\end{itemize}

% --- RB05 ---
\subsubsection{RB05 --- Ràng buộc kiểm soát hôn nhân và cấm cận huyết (Consanguinity Constraint)}
\begin{itemize}
    \item \textbf{Nội dung nghiệp vụ:} Hai cá nhân kết hôn phải khác giới tính và tuân thủ Luật Hôn nhân \& Gia đình Việt Nam: khoảng cách huyết thống tới Tổ tiên chung gần nhất (\textit{LCA -- Lowest Common Ancestor}) phải vượt quá phạm vi 3 đời.
    \item \textbf{Biểu diễn toán học:} Gọi $\text{Dist}(p_1, p_2)$ là khoảng cách huyết thống:
    \begin{alignat}{2}
        &\forall p_1, p_2 \in TV, \quad
        &&\text{KetHon}(p_1, p_2) \implies p_1.\text{GioiTinh} \ne p_2.\text{GioiTinh} \\[3pt]
        &\forall p_1, p_2 \in TV, \quad
        &&\text{KetHon}(p_1, p_2) \implies \text{Dist}(p_1, p_2) > 3
    \end{alignat}
    \item \textbf{Bối cảnh:} Quan hệ hôn phối, thực thể \texttt{THANH\_VIEN}, thuộc tính \texttt{GioiTinh}.
    \item \textbf{Quy tắc toàn vẹn:} Mọi đề nghị kết hôn phải thỏa mãn đồng thời điều kiện khác giới và khoảng cách huyết thống, ngược lại bị từ chối.
\end{itemize}

% ==============================================================================
% 1.4.2. NHÓM QUY TẮC KIỂM SOÁT TRUY CẬP VÀ PHÂN QUYỀN
% ==============================================================================
\subsection{Nhóm quy tắc kiểm soát truy cập và phân quyền sử dụng dữ liệu}

% --- RB06 ---
\subsubsection{RB06 --- Phân lập dữ liệu theo phạm vi Dòng họ (Family Data Isolation)}
\begin{itemize}
    \item \textbf{Nội dung nghiệp vụ:} Mỗi Dòng họ là một không gian thông tin riêng biệt và độc lập. Thành viên chỉ có quyền xem và thao tác dữ liệu (thành viên phả hệ, bài viết, sự kiện, di sản) thuộc dòng họ mình. Tư liệu được Trưởng tộc công bố công khai là ngoại lệ duy nhất được phép truy cập chéo dòng họ.
    \item \textbf{Biểu diễn toán học:} Gọi $\text{DHTV}$ là Dòng họ của thực thể. Đặt $\text{CK}$ là cờ công khai:
    \begin{equation}
        \begin{aligned}
            \forall u \in \text{ND},\ \forall obj \in \mathcal{E}: \quad
            \text{CoQuyenXem}(u, obj)
            &\iff \bigl[\,\text{DongHo}(u) = \text{DongHo}(obj)\,\bigr] \\
            &\lor\ \bigl[\,obj.\text{CongKhai} = \top\,\bigr]
        \end{aligned}
    \end{equation}
    \item \textbf{Bối cảnh:} Các thực thể \texttt{NGUOI\_DUNG}, \texttt{DONG\_HO}, \texttt{THANH\_VIEN}, \texttt{BAI\_VIET}, \texttt{SU\_KIEN}, \texttt{DI\_SAN\_LICH\_SU}.
    \item \textbf{Quy tắc toàn vẹn:} Mọi thao tác truy cập dữ liệu phải thỏa mãn $\text{CoQuyenXem}(u, obj) = \top$. Vi phạm điều kiện này là trái quy tắc nghiệp vụ và phải bị từ chối.
\end{itemize}

% --- RB07 ---
\subsubsection{RB07 --- Phân quyền chỉnh sửa hồ sơ theo vai trò Chi/Nhánh (Branch Role Authorization)}
\begin{itemize}
    \item \textbf{Nội dung nghiệp vụ:} Quyền chỉnh sửa hồ sơ thành viên phả hệ được phân cấp:
    \begin{itemize}[noitemsep]
        \item \textit{Trưởng chi họ:} Có quyền chỉnh sửa toàn bộ hồ sơ, bài viết và sự kiện trong chi nhánh mình quản lý; không can thiệp vào chi nhánh khác.
        \item \textit{Thành viên thường:} Chỉ chỉnh sửa hồ sơ bản thân và con ruột dưới 18 tuổi; có quyền xem toàn bộ cây phả hệ.
    \end{itemize}
    \item \textbf{Biểu diễn toán học:} Gọi $\text{ND} = \text{NGUOI\_DUNG}$, $TV = \text{THANH\_VIEN}$:
    \begin{equation}
        \begin{aligned}
            \text{CoQuyenChinhSua}(u, p) \iff\ &
            \bigl[\,u.\text{VaiTro} = \text{TruongChi}
            \land u.\text{MaChiToc} = p.\text{MaChiToc}\,\bigr] \\
            &\lor\; \bigl[\,u.\text{MaTV} = p.\text{MaTV}\,\bigr]
        \end{aligned}
    \end{equation}
    \item \textbf{Bối cảnh:} Thực thể \texttt{NGUOI\_DUNG}, \texttt{THANH\_VIEN}, \texttt{CHI\_TOC}, quan hệ phân quyền vai trò.
    \item \textbf{Quy tắc toàn vẹn:} Mọi thao tác cập nhật hồ sơ chỉ được thực hiện khi $\text{CoQuyenChinhSua}(u, p) = \top$. Không thành viên nào được sửa đổi hồ sơ ngoài phạm vi quyền hạn nghiệp vụ.
\end{itemize}

% --- RB08 ---
\subsubsection{RB08 --- Tính xác thực song ánh khi liên kết hồ sơ phả hệ (Profile Claiming 1:1)}
\begin{itemize}
    \item \textbf{Nội dung nghiệp vụ:} Một tài khoản người dùng chỉ được liên kết với tối đa một hồ sơ phả hệ và ngược lại, một hồ sơ phả hệ chỉ thuộc về tối đa một tài khoản đang kích hoạt (ánh xạ song ánh).
    \item \textbf{Biểu diễn toán học:}
    \begin{alignat}{2}
        &\forall u_1, u_2 \in \text{ND},\ \forall p \in TV, \quad
        &&\text{LienKet}(u_1, p) \land \text{LienKet}(u_2, p) \implies u_1 = u_2 \\[3pt]
        &\forall u \in \text{ND},\ \forall p_1, p_2 \in TV, \quad
        &&\text{LienKet}(u, p_1) \land \text{LienKet}(u, p_2) \implies p_1 = p_2
    \end{alignat}
    \item \textbf{Bối cảnh:} Thực thể \texttt{NGUOI\_DUNG}, \texttt{THANH\_VIEN}, quan hệ liên kết tài khoản -- hồ sơ.
    \item \textbf{Quy tắc toàn vẹn:} Yêu cầu liên kết hồ sơ phải bị từ chối nếu thành viên hoặc tài khoản đã có liên kết hợp lệ trước đó.
\end{itemize}

% --- RB09 ---
\subsubsection{RB09 --- Tính nhất quán mốc thời gian sự kiện (Event Timeline Consistency)}
\begin{itemize}
    \item \textbf{Nội dung nghiệp vụ:} Thời điểm kết thúc của một sự kiện dòng họ (nếu có ghi nhận) bắt buộc phải diễn ra sau thời điểm bắt đầu.
    \item \textbf{Biểu diễn toán học:} Gọi $SK = \text{SU\_KIEN}$, $t_s = \text{ThoiGianBatDau}$, $t_e = \text{ThoiGianKetThuc}$:
    \begin{equation}
        \forall e \in SK: \quad
        e.t_e \ne \emptyset \implies e.t_e > e.t_s
    \end{equation}
    \item \textbf{Bối cảnh:} Thực thể \texttt{SU\_KIEN}, thuộc tính \texttt{ThoiGianBatDau}, \texttt{ThoiGianKetThuc}.
    \item \textbf{Quy tắc toàn vẹn:} Mọi bản ghi sự kiện có thời gian kết thúc đều phải thỏa mãn $t_e > t_s$. Dữ liệu vi phạm không có ý nghĩa trong thực tế và không được phép tồn tại.
\end{itemize}

% --- RB10 ---
\subsubsection{RB10 --- Tính nhất quán trạng thái sinh tử và ngày mất (Vital Status Consistency)}
\begin{itemize}
    \item \textbf{Nội dung nghiệp vụ:} Nếu đã ghi nhận ngày mất của một thành viên thì trạng thái sinh học bắt buộc phải là ``Đã mất'' và ngày mất không thể xảy ra trước ngày sinh.
    \item \textbf{Biểu diễn toán học:} Gọi $\text{nm} = \text{NgayMat}$, $\text{ns} = \text{NgaySinh}$, $\text{tt} = \text{TinhTrang}$:
    \begin{equation}
        \begin{aligned}
            \forall p \in TV: \quad p.\text{nm} \ne \emptyset \implies
            &\;\bigl(\,p.\text{tt} = \text{Đã Mất}\,\bigr) \\
            &\land\; \bigl(\,p.\text{nm} \ge p.\text{ns}\,\bigr)
        \end{aligned}
    \end{equation}
    \item \textbf{Bối cảnh:} Thực thể \texttt{THANH\_VIEN}, thuộc tính \texttt{TinhTrang}, \texttt{NgayMat}, \texttt{NgaySinhDuongLich}.
    \item \textbf{Quy tắc toàn vẹn:} Không thể đồng thời tồn tại hồ sơ có ngày mất được ghi nhận mà trạng thái vẫn là ``Còn sống'', hoặc có ngày mất trước ngày sinh. Hai trạng thái này mâu thuẫn với thực tế sinh học và phải bị ràng buộc loại trừ.
\end{itemize}

% ==============================================================================
% 1.4.3. BẢNG TẦM ẢNH HƯỞNG CỦA CÁC RÀNG BUỘC TOÀN VẸN
% ==============================================================================
\subsection{Bảng tầm ảnh hưởng của các ràng buộc toàn vẹn (Impact Matrix)}
Bảng tầm ảnh hưởng xác định các thao tác biến đổi dữ liệu: \textbf{Thêm mới} ($+$), \textbf{Xóa} ($-$), \textbf{Sửa đổi} ($\ast$) trên từng nhóm thực thể có nguy cơ vi phạm ràng buộc toàn vẹn. Kết quả phân tích này là đầu vào cho giai đoạn thiết kế logic và vật lý:

\begin{table}[H]
\centering
\small
\setlength{\tabcolsep}{3pt}
\renewcommand{\arraystretch}{1.25}
\begin{tabularx}{\textwidth}{|l|*{12}{>{\centering\arraybackslash}X|}}
\hline
\rowcolor{tableheader}
\textcolor{white}{\textbf{Mã RBTV}} &
\multicolumn{3}{c|}{\textcolor{white}{\textbf{THANH\_VIEN}}} &
\multicolumn{3}{c|}{\textcolor{white}{\textbf{NGUOI\_DUNG}}} &
\multicolumn{3}{c|}{\textcolor{white}{\textbf{BAI\_VIET / EVT}}} &
\multicolumn{3}{c|}{\textcolor{white}{\textbf{DONG\_HO / CHI}}} \\ \cline{2-13}
\rowcolor{tableheader}
& \textcolor{white}{\small\textbf{+}} & \textcolor{white}{\small\textbf{$-$}} & \textcolor{white}{\small\textbf{$\ast$}}
& \textcolor{white}{\small\textbf{+}} & \textcolor{white}{\small\textbf{$-$}} & \textcolor{white}{\small\textbf{$\ast$}}
& \textcolor{white}{\small\textbf{+}} & \textcolor{white}{\small\textbf{$-$}} & \textcolor{white}{\small\textbf{$\ast$}}
& \textcolor{white}{\small\textbf{+}} & \textcolor{white}{\small\textbf{$-$}} & \textcolor{white}{\small\textbf{$\ast$}} \\ \hline
\rowcolor{tablerowalt}
\textbf{RB01} & $\checkmark$ & & $\checkmark$ & & & & & & & & & \\ \hline
\textbf{RB02} & $\checkmark$ & & $\checkmark$ & & & & & & & & & \\ \hline
\rowcolor{tablerowalt}
\textbf{RB03} & $\checkmark$ & & $\checkmark$ & & & & & & & & & \\ \hline
\textbf{RB04} & $\checkmark$ & & $\checkmark$ & & & & & & & & & \\ \hline
\rowcolor{tablerowalt}
\textbf{RB05} & $\checkmark$ & & $\checkmark$ & & & & & & & & & \\ \hline
\textbf{RB06} & $\checkmark$ & & $\checkmark$ & $\checkmark$ & & $\checkmark$ & $\checkmark$ & & $\checkmark$ & $\checkmark$ & $\checkmark$ & $\checkmark$ \\ \hline
\rowcolor{tablerowalt}
\textbf{RB07} & $\checkmark$ & $\checkmark$ & $\checkmark$ & & & $\checkmark$ & $\checkmark$ & $\checkmark$ & $\checkmark$ & & & \\ \hline
\textbf{RB08} & & & $\checkmark$ & $\checkmark$ & & $\checkmark$ & & & & & & \\ \hline
\rowcolor{tablerowalt}
\textbf{RB09} & & & & & & & $\checkmark$ & & $\checkmark$ & & & \\ \hline
\textbf{RB10} & $\checkmark$ & & $\checkmark$ & & & & & & & & & \\ \hline
\end{tabularx}
\caption{Bảng tầm ảnh hưởng của các ràng buộc toàn vẹn mức quan niệm (Impact Matrix)}
\label{tab:impact_matrix_conceptual}
\end{table}

\noindent\textbf{Quy ước ký hiệu:}\ Dấu $(\checkmark)$ chỉ thao tác có nguy cơ vi phạm ràng buộc và bắt buộc phải kiểm tra tính hợp lệ. Ô trống có nghĩa thao tác tương ứng không ảnh hưởng đến ràng buộc này.
\begin{itemize}[noitemsep,topsep=4pt]
    \item \textbf{Cột $+$ (Thêm mới):} Ghi nhận bản ghi mới có nguy cơ vi phạm ràng buộc.
    \item \textbf{Cột $-$ (Xóa):} Xóa bản ghi có thể phá vỡ tham chiếu toàn vẹn.
    \item \textbf{Cột $\ast$ (Sửa đổi):} Cập nhật thuộc tính liên quan có nguy cơ vi phạm ràng buộc.
\end{itemize}



% Chương 2: Thiết kế Logic (Logical Design)
% !TEX root = ../../main.tex
\chapter{Thiết kế Logic (Logical Design)}
\label{chap:chuong2_logic}

Giai đoạn thiết kế logic thực hiện chuyển đổi mô hình Thực thể -- Kết hợp trừu tượng thành Lược đồ quan hệ (Relational Schema) tương thích với mô hình dữ liệu quan hệ, tiến hành chuẩn hóa dữ liệu để triệt tiêu dị thường và xác lập các ràng buộc toàn vẹn logic.

% !TEX root = ../../main.tex
\section{Quy tắc chuyển đổi từ mô hình ER sang mô hình quan hệ (Relational Model)}
\label{sec:2_1_quy_tac_chuyen_doi}

Chuyển đổi từ mô hình Thực thể -- Liên kết (ERD) sang Mô hình Quan hệ (Relational Model) là bước then chốt trong giai đoạn Thiết kế Logic (Logical Design). Quá trình này đảm bảo toàn bộ thông tin, ràng buộc dữ liệu và mối quan hệ giữa các đối tượng trong hệ thống Quản lý Gia phả được biểu diễn chính xác dưới dạng các lược đồ quan hệ (Relation schemas).

\subsection{Chuyển đổi thực thể thông thường thành các bảng (Relation schemas)}
\label{subsec:2_1_1_chuyen_doi_thuc_the}

Mỗi thực thể mạnh trong mô hình ERD tương ứng với một bảng trong mô hình quan hệ. Các thuộc tính đơn trở thành các cột của bảng, và thuộc tính định danh được chọn làm Khóa chính (Primary Key -- PK).

\begin{itemize}
	\item \textbf{Thực thể Thành viên (\texttt{ThanhVien}):} Lưu trữ thông tin định danh và cá nhân của từng người trong gia tộc.
	\begin{itemize}
		\item Lược đồ quan hệ: \texttt{THANH\_VIEN}(\underline{\texttt{MaThanhVien}}, \texttt{HoTen}, \texttt{GioiTinh}, \texttt{NgaySinh}, \texttt{NgayMat}, \texttt{DiaChi}, \texttt{TrangThaiConSong}).
		\item \textbf{PK:} \texttt{MaThanhVien}.
	\end{itemize}
	
	\item \textbf{Thực thể Bài viết (\texttt{BaiViet}):} Lưu trữ tin tức, bài viết truyền thống của gia tộc.
	\begin{itemize}
		\item Lược đồ quan hệ: \texttt{BAI\_VIET}(\underline{\texttt{MaBaiViet}}, \texttt{TieuDe}, \texttt{NoiDung}, \texttt{ThoiGianDang}, \texttt{QuyenRiengTu}).
		\item \textbf{PK:} \texttt{MaBaiViet}.
	\end{itemize}
	
	\item \textbf{Thực thể Sự kiện (\texttt{SuKien}):} Lưu trữ thông tin ngày giỗ, họp mặt, giỗ tổ.
	\begin{itemize}
		\item Lược đồ quan hệ: \texttt{SU\_KIEN}(\underline{\texttt{MaSuKien}}, \texttt{TenSuKien}, \texttt{NgayAmLich}, \texttt{NgayDuongLich}, \texttt{DiaDiem}).
		\item \textbf{PK:} \texttt{MaSuKien}.
	\end{itemize}
\end{itemize}

\subsection{Chuyển đổi các mối quan hệ 1--1, 1--N, N--N và mối quan hệ đệ quy}
\label{subsec:2_1_2_chuyen_doi_moi_quan_he}

\subsubsection{Mối quan hệ 1--1 (One-to-One)}
Được xử lý bằng cách đưa Khóa chính của một trong hai thực thể sang thực thể còn lại làm Khóa ngoại (Foreign Key -- FK).
\begin{itemize}
	\item \textbf{Tài khoản -- Thành viên (1--1):} Mỗi thành viên có thể được cấp một tài khoản đăng nhập hệ thống.
	\item Lược đồ quan hệ: \texttt{TAI\_KHOAN}(\underline{\texttt{MaTaiKhoan}}, \texttt{TenDangNhap}, \texttt{MatKhau}, \texttt{VaiTro}, \texttt{MaThanhVien}\textsuperscript{FK}).
\end{itemize}

\subsubsection{Mối quan hệ 1--N (One-to-Many)}
Khóa chính của bảng phía nhánh 1 được đưa sang làm Khóa ngoại ở bảng phía nhánh N.
\begin{itemize}
	\item \textbf{Tác giả bài viết (Thành viên -- Bài viết):} Một thành viên có thể đăng nhiều bài viết.
	\item Lược đồ quan hệ: \texttt{BAI\_VIET}(\underline{\texttt{MaBaiViet}}, \texttt{TieuDe}, \texttt{NoiDung}, \texttt{ThoiGianDang}, \texttt{MaNguoiDang}\textsuperscript{FK}).
\end{itemize}

\subsubsection{Mối quan hệ N--N (Many-to-Many)}
Tạo một bảng liên kết trung gian chứa Khóa chính của hai thực thể tham gia làm Khóa ngoại. Tập hợp hai khóa ngoại này hợp thành Khóa chính hợp phần.
\begin{itemize}
	\item \textbf{Tham gia sự kiện (Thành viên -- Sự kiện):}
	\item Lược đồ quan hệ: \texttt{THAM\_GIA\_SU\_KIEN}(\underline{\texttt{MaThanhVien}}\textsuperscript{FK}, \underline{\texttt{MaSuKien}}\textsuperscript{FK}, \texttt{TrangThaiRSVP}, \texttt{GhiChu}).
\end{itemize}

\subsubsection{Mối quan hệ đệ quy (Recursive Relationships)}
Áp dụng cho các mối quan hệ nội bộ giữa các đối tượng trong cùng thực thể \texttt{ThanhVien}:
\begin{itemize}
	\item \textbf{Mối quan hệ Cha/Mẹ -- Con (Đệ quy 1--N):} Mỗi thành viên có một cha và một mẹ (đều thuộc bảng \texttt{THANH\_VIEN}). Ta thêm hai thuộc tính khóa ngoại tự tham chiếu:
	\begin{itemize}
		\item Lược đồ quan hệ bổ sung: \texttt{THANH\_VIEN}(\dots, \texttt{MaCha}\textsuperscript{FK}, \texttt{MaMe}\textsuperscript{FK}).
	\end{itemize}
	\item \textbf{Mối quan hệ Hôn phối / Vợ -- Chồng (Đệ quy 1--1 hoặc N--N theo thời gian):} Để lưu lịch sử hôn nhân (kết hôn, ly hôn), tạo bảng liên kết đệ quy:
	\begin{itemize}
		\item Lược đồ quan hệ: \texttt{QUAN\_HE\_HON\_NHAN}(\underline{\texttt{MaChong}}\textsuperscript{FK}, \underline{\texttt{MaVo}}\textsuperscript{FK}, \texttt{NgayKetHon}, \texttt{TrangThai}).
	\end{itemize}
\end{itemize}

\subsection{Xử lý quan hệ kế thừa hoặc phân loại}
\label{subsec:2_1_3_xu_ly_ke_thua}

Trong cây gia phả, thực thể \texttt{ThanhVien} được phân loại (kế thừa) thành hai nhóm: \textbf{Thành viên nội tộc (Ruột thịt)} và \textbf{Thành viên ngoại tộc (Dâu/Rể)}. 

Áp dụng phương pháp \textbf{Biểu diễn lớp cha và các lớp con (Class Hierarchy Mapping)} để tránh dư thừa dữ liệu:

\begin{enumerate}
	\item \textbf{Bảng lớp cha (\texttt{THANH\_VIEN}):} Chứa tất cả các thuộc tính chung (Họ tên, Ngày sinh, Giới tính, \dots).
	\item \textbf{Bảng lớp con \texttt{THANH\_VIEN\_NOI\_TOC}:} Chứa thuộc tính riêng của người mang dòng máu gia tộc (Đời thứ mấy, Chi/Nhánh). Khóa chính \texttt{MaThanhVien} vừa là PK vừa là FK tham chiếu về \texttt{THANH\_VIEN}.
	\item \textbf{Bảng lớp con \texttt{THANH\_VIEN\_DAU\_RE}:} Chứa thuộc tính riêng của dâu/rể (Gia tộc gốc, Ngày nhập tộc). Khóa chính \texttt{MaThanhVien} vừa là PK vừa là FK tham chiếu về \texttt{THANH\_VIEN}.
\end{enumerate}

\textbf{Lược đồ chi tiết:}
\begin{itemize}
	\item \texttt{THANH\_VIEN\_NOI\_TOC}(\underline{\texttt{MaThanhVien}}\textsuperscript{PK, FK}, \texttt{DoiThu}, \texttt{ThuocChi}).
	\item \texttt{THANH\_VIEN\_DAU\_RE}(\underline{\texttt{MaThanhVien}}\textsuperscript{PK, FK}, \texttt{GiaTocGoc}, \texttt{NgayVeLamDauRe}).
\end{itemize}

% !TEX root = ../../main.tex
\section{Lược đồ quan hệ logic (Relational Schemas)}
\label{sec:2_2_luoc_do_quan_he}

Sau khi áp dụng các quy tắc chuyển đổi từ mô hình thực thể -- liên kết (ERD), hệ thống cơ sở dữ liệu \textbf{FamilyConnect} được cụ thể hóa thành tập hợp các lược đồ quan hệ (Relation Schemas). Dưới đây là mô tả chi tiết các lược đồ quan hệ thu được.

\subsection{Lược đồ các bảng quan hệ thực thể}
\label{subsec:2_2_1_chuyen_doi_thuc_the}

Mỗi thực thể mạnh trong mô hình ERD tương ứng với một bảng trong mô hình quan hệ. Các thuộc tính đơn trở thành các cột của bảng, và thuộc tính định danh được chọn làm Khóa chính (Primary Key -- PK).

\begin{itemize}
	\item \textbf{Thực thể Thành viên (\texttt{ThanhVien}):} Lưu trữ thông tin định danh và cá nhân của từng người trong gia tộc.
	\begin{itemize}
		\item Lược đồ quan hệ: \texttt{THANH\_VIEN}(\underline{\texttt{MaThanhVien}}, \texttt{HoTen}, \texttt{GioiTinh}, \texttt{NgaySinh}, \texttt{NgayMat}, \texttt{DiaChi}, \texttt{TrangThaiConSong}).
		\item \textbf{PK:} \texttt{MaThanhVien}.
	\end{itemize}
	
	\item \textbf{Thực thể Bài viết (\texttt{BaiViet}):} Lưu trữ tin tức, bài viết truyền thống của gia tộc.
	\begin{itemize}
		\item Lược đồ quan hệ: \texttt{BAI\_VIET}(\underline{\texttt{MaBaiViet}}, \texttt{TieuDe}, \texttt{NoiDung}, \texttt{ThoiGianDang}, \texttt{QuyenRiengTu}).
		\item \textbf{PK:} \texttt{MaBaiViet}.
	\end{itemize}
	
	\item \textbf{Thực thể Sự kiện (\texttt{SuKien}):} Lưu trữ thông tin ngày giỗ, họp mặt, giỗ tổ.
	\begin{itemize}
		\item Lược đồ quan hệ: \texttt{SU\_KIEN}(\underline{\texttt{MaSuKien}}, \texttt{TenSuKien}, \texttt{NgayAmLich}, \texttt{NgayDuongLich}, \texttt{DiaDiem}).
		\item \textbf{PK:} \texttt{MaSuKien}.
	\end{itemize}
\end{itemize}

\subsection{Chuyển đổi các mối quan hệ 1--1, 1--N, N--N và mối quan hệ đệ quy}
\label{subsec:2_2_2_chuyen_doi_moi_quan_he}

\subsubsection{Mối quan hệ 1--1 (One-to-One)}
Được xử lý bằng cách đưa Khóa chính của một trong hai thực thể sang thực thể còn lại làm Khóa ngoại (Foreign Key -- FK).
\begin{itemize}
	\item \textbf{Tài khoản -- Thành viên (1--1):} Mỗi thành viên có thể được cấp một tài khoản đăng nhập hệ thống.
	\item Lược đồ quan hệ: \texttt{TAI\_KHOAN}(\underline{\texttt{MaTaiKhoan}}, \texttt{TenDangNhap}, \texttt{MatKhau}, \texttt{VaiTro}, \texttt{MaThanhVien}\textsuperscript{FK}).
\end{itemize}

\subsubsection{Mối quan hệ 1--N (One-to-Many)}
Khóa chính của bảng phía nhánh 1 được đưa sang làm Khóa ngoại ở bảng phía nhánh N.
\begin{itemize}
	\item \textbf{Tác giả bài viết (Thành viên -- Bài viết):} Một thành viên có thể đăng nhiều bài viết.
	\item Lược đồ quan hệ: \texttt{BAI\_VIET}(\underline{\texttt{MaBaiViet}}, \texttt{TieuDe}, \texttt{NoiDung}, \texttt{ThoiGianDang}, \texttt{MaNguoiDang}\textsuperscript{FK}).
\end{itemize}

\subsubsection{Mối quan hệ N--N (Many-to-Many)}
Tạo một bảng liên kết trung gian chứa Khóa chính của hai thực thể tham gia làm Khóa ngoại. Tập hợp hai khóa ngoại này hợp thành Khóa chính hợp phần.
\begin{itemize}
	\item \textbf{Tham gia sự kiện (Thành viên -- Sự kiện):}
	\item Lược đồ quan hệ: \texttt{THAM\_GIA\_SU\_KIEN}(\underline{\texttt{MaThanhVien}}\textsuperscript{FK}, \underline{\texttt{MaSuKien}}\textsuperscript{FK}, \texttt{TrangThaiRSVP}, \texttt{GhiChu}).
\end{itemize}

\subsubsection{Mối quan hệ đệ quy (Recursive Relationships)}
Áp dụng cho các mối quan hệ nội bộ giữa các đối tượng trong cùng thực thể \texttt{ThanhVien}:
\begin{itemize}
	\item \textbf{Mối quan hệ Cha/Mẹ -- Con (Đệ quy 1--N):} Mỗi thành viên có một cha và một mẹ (đều thuộc bảng \texttt{THANH\_VIEN}). Ta thêm hai thuộc tính khóa ngoại tự tham chiếu:
	\begin{itemize}
		\item Lược đồ quan hệ bổ sung: \texttt{THANH\_VIEN}(\dots, \texttt{MaCha}\textsuperscript{FK}, \texttt{MaMe}\textsuperscript{FK}).
	\end{itemize}
	\item \textbf{Mối quan hệ Hôn phối / Vợ -- Chồng (Đệ quy 1--1 hoặc N--N theo thời gian):} Để lưu lịch sử hôn nhân (kết hôn, ly hôn), tạo bảng liên kết đệ quy:
	\begin{itemize}
		\item Lược đồ quan hệ: \texttt{QUAN\_HE\_HON\_NHAN}(\underline{\texttt{MaChong}}\textsuperscript{FK}, \underline{\texttt{MaVo}}\textsuperscript{FK}, \texttt{NgayKetHon}, \texttt{TrangThai}).
	\end{itemize}
\end{itemize}

\subsection{Xử lý quan hệ kế thừa hoặc phân loại}
\label{subsec:2_2_3_xu_ly_ke_thua}

Trong cây gia phả, thực thể \texttt{ThanhVien} được phân loại (kế thừa) thành hai nhóm: \textbf{Thành viên nội tộc (Ruột thịt)} và \textbf{Thành viên ngoại tộc (Dâu/Rể)}. 

Áp dụng phương pháp \textbf{Biểu diễn lớp cha và các lớp con (Class Hierarchy Mapping)} để tránh dư thừa dữ liệu:

\begin{enumerate}
	\item \textbf{Bảng lớp cha (\texttt{THANH\_VIEN}):} Chứa tất cả các thuộc tính chung (Họ tên, Ngày sinh, Giới tính, \dots).
	\item \textbf{Bảng lớp con \texttt{THANH\_VIEN\_NOI\_TOC}:} Chứa thuộc tính riêng của người mang dòng máu gia tộc (Đời thứ mấy, Chi/Nhánh). Khóa chính \texttt{MaThanhVien} vừa là PK vừa là FK tham chiếu về \texttt{THANH\_VIEN}.
	\item \textbf{Bảng lớp con \texttt{THANH\_VIEN\_DAU\_RE}:} Chứa thuộc tính riêng của dâu/rể (Gia tộc gốc, Ngày nhập tộc). Khóa chính \texttt{MaThanhVien} vừa là PK vừa là FK tham chiếu về \texttt{THANH\_VIEN}.
\end{enumerate}

\textbf{Lược đồ chi tiết:}
\begin{itemize}
	\item \texttt{THANH\_VIEN\_NOI\_TOC}(\underline{\texttt{MaThanhVien}}\textsuperscript{PK, FK}, \texttt{DoiThu}, \texttt{ThuocChi}).
	\item \texttt{THANH\_VIEN\_DAU\_RE}(\underline{\texttt{MaThanhVien}}\textsuperscript{PK, FK}, \texttt{GiaTocGoc}, \texttt{NgayVeLamDauRe}).
\end{itemize}

% !TEX root = ../../main.tex
\section{Chuẩn hóa dữ liệu (Normalization)}
\label{sec:2_3_chuan_hoa_du_lieu}

Chuẩn hóa là quá trình tổ chức lại lược đồ quan hệ nhằm loại bỏ sự dư thừa dữ liệu và các bất thường cập nhật (thêm, sửa, xóa). Trong phần này, chúng ta áp dụng lý thuyết \textbf{Phụ thuộc hàm (Functional Dependency -- FD)} để phân tích 3 quan hệ trọng yếu của hệ thống FamilyConnect, xác định khóa ứng viên và đánh giá mức độ chuẩn hóa.

\subsection{Phân tích Phụ thuộc hàm (Functional Dependencies)}
\label{subsec:2_3_1_phu_thuoc_ham}

\subsubsection{Quan hệ MemberProfile}

\begin{example}[Phân tích FD cho MemberProfile trước khi phân rã]
Xét lược đồ \emph{trước khi tách} quan hệ \texttt{Branch}: \\
$MemberProfile\_raw(M, U, E, B, BN, G)$, trong đó:
\begin{itemize}[noitemsep]
    \item $M$ = member\_id, $U$ = user\_id, $E$ = email (từ Users)
    \item $B$ = branch\_id, $BN$ = branch\_name, $G$ = generation\_level
\end{itemize}
Tập phụ thuộc hàm $F_1$:
\begin{align*}
    f_1: & \quad M \rightarrow \{U, B, BN, G\} \\
    f_2: & \quad U \rightarrow E \\
    f_3: & \quad B \rightarrow \{BN, G\}
\end{align*}

\textbf{Tính bao đóng khóa:} $\{M\}^+ = \{M\} \xrightarrow{f_1} \{M,U,B,BN,G\} \xrightarrow{f_2} \{M,U,B,BN,G,E\} = U_{\text{all}}$. \\
Vậy $K_1 = \{M\}$ là \textbf{Khóa ứng viên duy nhất}.

\textbf{Vi phạm BCNF:} $f_3: B \rightarrow \{BN, G\}$ nhưng $B$ không phải siêu khóa. \\
$\implies$ Tồn tại phụ thuộc bắc cầu $M \rightarrow B \rightarrow \{BN,G\}$, vi phạm \textbf{3NF và BCNF}.
\end{example}

\subsubsection{Quan hệ Users}

\begin{example}[Phân tích FD cho Users]
$Users(U, UN, EM, PH, FN, CA, IA, RC)$, trong đó $U$ = user\_id, $UN$ = username, $EM$ = email, $PH$ = phone\_number, $FN$ = full\_name, $CA$ = created\_at, $IA$ = is\_active, $RC$ = role\_code.

Tập phụ thuộc hàm $F_2$:
\begin{align*}
    g_1: & \quad U \rightarrow \{UN, EM, PH, FN, CA, IA, RC\} \\
    g_2: & \quad UN \rightarrow U \quad (\text{unique constraint}) \\
    g_3: & \quad EM \rightarrow U \quad (\text{unique constraint})
\end{align*}

\textbf{Khóa ứng viên:} $K = \{U\}$, $K' = \{UN\}$, $K'' = \{EM\}$. \\
Vì mọi phụ thuộc hàm đều xuất phát từ siêu khóa, lược đồ \textbf{đạt BCNF}.
\end{example}

\subsubsection{Quan hệ MarriageRelation}

\begin{example}[Phân tích FD cho MarriageRelation]
$MarriageRelation(MID, HID, WID, MD, DD, ST)$, trong đó $MID$ = marriage\_id, $HID$ = husband\_id, $WID$ = wife\_id, $MD$ = marriage\_date, $DD$ = divorce\_date, $ST$ = status.

Tập phụ thuộc hàm $F_3$:
\begin{align*}
    h_1: & \quad MID \rightarrow \{HID, WID, MD, DD, ST\} \\
    h_2: & \quad (HID, WID, MD) \rightarrow MID
\end{align*}

\textbf{Khóa ứng viên:} $K = \{MID\}$ và $K' = \{HID, WID, MD\}$. \\
Không có phụ thuộc một phần hay bắc cầu, lược đồ \textbf{đạt BCNF}.
\end{example}

\subsection{Đánh giá các Dạng chuẩn (1NF, 2NF, 3NF, BCNF)}
\label{subsec:2_3_2_danh_gia_dang_chuan}

\subsubsection{Tiêu chí từng dạng chuẩn}

\begin{table}[htbp]
\centering
\small
\caption{Tiêu chí đánh giá các dạng chuẩn.}
\label{tab:dang_chuan_tieu_chi}
\begin{tabularx}{\textwidth}{l Y Y}
\toprule
\textbf{Dạng chuẩn} & \textbf{Yêu cầu} & \textbf{Loại bỏ} \\
\midrule
1NF & Mọi thuộc tính là nguyên tử, không có nhóm lặp & Thuộc tính đa trị / hợp thành \\
2NF & Đạt 1NF + không có phụ thuộc hàm một phần vào khóa & Phụ thuộc một phần (partial FD) \\
3NF & Đạt 2NF + không có phụ thuộc bắc cầu qua thuộc tính không khóa & Phụ thuộc bắc cầu (transitive FD) \\
BCNF & Mọi vế trái của FD đều phải là siêu khóa & Vi phạm BCNF còn sót sau 3NF \\
\bottomrule
\end{tabularx}
\end{table}

\subsubsection{Kết quả đánh giá các lược đồ FamilyConnect}

\begin{table}[htbp]
\centering
\small
\caption{Đánh giá mức độ chuẩn hóa các lược đồ quan hệ chính của FamilyConnect.}
\label{tab:so_sanh_dang_chuan}
\begin{tabularx}{\textwidth}{l c c c c Y}
\toprule
\textbf{Lược đồ} & \textbf{1NF} & \textbf{2NF} & \textbf{3NF} & \textbf{BCNF} & \textbf{Ghi chú} \\
\midrule
Users                & \checkmark & \checkmark & \checkmark & \checkmark & Đạt BCNF \\
FamilyTree           & \checkmark & \checkmark & \checkmark & \checkmark & Đạt BCNF \\
Branch               & \checkmark & \checkmark & \checkmark & \checkmark & Đạt BCNF \\
MemberProfile (sau)  & \checkmark & \checkmark & \checkmark & \checkmark & Đạt BCNF sau phân rã \\
ParentChildRelation  & \checkmark & \checkmark & \checkmark & \checkmark & Khóa hợp phần, đạt BCNF \\
MarriageRelation     & \checkmark & \checkmark & \checkmark & \checkmark & 2 khóa ứng viên, đạt BCNF \\
Post                 & \checkmark & \checkmark & \checkmark & \checkmark & Đạt BCNF \\
PostInteraction      & \checkmark & \checkmark & \checkmark & \checkmark & Đạt BCNF \\
Event                & \checkmark & \checkmark & \checkmark & \checkmark & Đạt BCNF \\
EventRSVP            & \checkmark & \checkmark & \checkmark & \checkmark & Đạt BCNF \\
HeritageDocument     & \checkmark & \checkmark & \checkmark & \checkmark & Đạt BCNF \\
AIAssistantLog       & \checkmark & \checkmark & \checkmark & \checkmark & Đạt BCNF \\
AuditLog             & \checkmark & \checkmark & \checkmark & \checkmark & Đạt BCNF \\
\bottomrule
\end{tabularx}
\end{table}

\subsubsection{Phân rã lược đồ MemberProfile đạt BCNF}

Lược đồ \texttt{MemberProfile\_raw} vi phạm BCNF do $f_3: B \rightarrow \{BN,G\}$ với $B$ không là siêu khóa. Áp dụng \textbf{thuật toán phân rã BCNF (lossless-join decomposition)}:

\begin{enumerate}[label=\arabic*.]
    \item Tách phụ thuộc vi phạm $f_3$ ra thành lược đồ con:
    \[
    \textbf{Branch}(\underline{branch\_id}, branch\_name, generation\_level)
    \quad \text{với } F_B = \{B \rightarrow \{BN, G\}\} \quad \text{đạt BCNF}
    \]
    \item Lược đồ còn lại sau khi tách:
    \[
    \textbf{MemberProfile}(\underline{member\_id}, user\_id, branch\_id^{FK}, \ldots)
    \quad \text{với } F_M = \{M \rightarrow \{U,B\},\; U \rightarrow E\} \quad \text{đạt BCNF}
    \]
\end{enumerate}

\textbf{Kiểm tra tính chất bảo toàn:}
\begin{itemize}[noitemsep]
    \item \textbf{Lossless-join:} $Branch \bowtie MemberProfile$ khôi phục đầy đủ quan hệ gốc vì $Branch \cap MemberProfile = \{branch\_id\}$ là siêu khóa của Branch. $\checkmark$
    \item \textbf{Dependency preservation:} $F_B \cup F_M$ bao phủ đầy đủ $F_1$. $\checkmark$
\end{itemize}

% !TEX root = ../../main.tex
\section{Ràng buộc Toàn vẹn Mức Logic (Integrity Constraints)}
\label{sec:2_4_rang_buoc_logic}

Ràng buộc toàn vẹn mức logic là các quy tắc nghiệp vụ bắt buộc hệ quản trị CSDL phải tuân thủ nhằm đảm bảo tính đúng đắn, nhất quán của dữ liệu. Hệ thống \textbf{FamilyConnect} áp dụng 4 nhóm ràng buộc toàn vẹn chính.

\subsection{Ràng buộc Miền giá trị (Domain Constraints)}
\label{subsec:2_4_1_domain}

\begin{itemize}
    \item \textbf{C1.1 (Gender Domain):} Thuộc tính $\text{gender} \in \{\text{'Male'}, \text{'Female'}, \text{'Other'}\}$.
    \item \textbf{C1.2 (RSVP Status Domain):} Thuộc tính $\text{rsvp\_status} \in \{\text{'Attending'}, \text{'Declined'}, \text{'Pending'}, \text{'Maybe'}\}$.
    \item \textbf{C1.3 (Reaction Type Domain):} Thuộc tính $\text{reaction\_type} \in \{\text{'Like'}, \text{'Love'}, \text{'Celebrate'}, \text{'Care'}\}$.
\end{itemize}

\subsection{Ràng buộc Thực thể và Khóa (Entity \& Key Constraints)}
\label{subsec:2_4_2_entity_key}

\begin{itemize}
    \item \textbf{C2.1 (Primary Keys):} Tất cả các bảng phải có Khóa chính (Primary Key) xác định duy nhất từng bản ghi và không được phép nhận giá trị \texttt{NULL}.
    \item \textbf{C2.2 (User Unique Identity):} Thuộc tính \texttt{username} và \texttt{email} trong bảng \texttt{Users} phải là duy nhất (\texttt{UNIQUE}).
    \item \textbf{C2.3 (Marriage Uniqueness):} Cặp $(\text{husband\_id}, \text{wife\_id}, \text{marriage\_date})$ trong quan hệ \texttt{MarriageRelation} là duy nhất.
\end{itemize}

\subsection{Ràng buộc Toàn vẹn Tham chiếu (Referential Integrity Constraints)}
\label{subsec:2_4_3_referential}

\begin{itemize}
    \item \textbf{C3.1 (Member to User Reference):}
    \[
    \forall m \in \text{MemberProfile}, \quad (m.\text{user\_id} \neq \text{NULL}) \implies (\exists u \in \text{Users}: m.\text{user\_id} = u.\text{user\_id})
    \]
    Chính sách hành vi khi xóa người dùng: \texttt{ON DELETE SET NULL}.

    \item \textbf{C3.2 (Parent-Child Reference):}
    \[
    \begin{aligned}
    \forall p \in \text{ParentChildRelation}, \ \exists m_1, m_2 \in \text{MemberProfile} \text{ sao cho: } \\
    (p.\text{parent\_id} = m_1.\text{member\_id}) \land (p.\text{child\_id} = m_2.\text{member\_id})
    \end{aligned}
    \]
    Chính sách hành vi khi xóa thành viên: \texttt{ON DELETE CASCADE}.
\end{itemize}

\subsection{Ràng buộc Nghiệp vụ Phức tạp (Cross-Table \& Semantic Constraints)}
\label{subsec:2_4_4_semantic}

\begin{definition}[Ràng buộc C4.1: Ngày sinh -- Ngày mất và Vòng đời]
Với mọi thành viên $m \in \text{MemberProfile}$:
\[
\begin{aligned}
m.\text{is\_alive} = \text{\texttt{FALSE}} &\implies m.\text{death\_date} \ge m.\text{birth\_date} \\
m.\text{is\_alive} = \text{\texttt{TRUE}}  &\implies m.\text{death\_date} \text{ IS NULL}
\end{aligned}
\]
\end{definition}

\begin{definition}[Ràng buộc C4.2: Ngăn ngừa Vòng lặp Quan hệ Huyết thống (Acyclic Genealogic Constraint)]
Đồ thị quan hệ cha mẹ -- con cái $G = (V, E)$ trong đó $V = \{\text{member\_id}\}$ và $E = \{(\text{parent\_id}, \text{child\_id})\}$ bắt buộc phải là một \textbf{Đồ thị có hướng không chu trình (Directed Acyclic Graph -- DAG)}.
\[
\forall m \in \text{MemberProfile}, \quad (m, m) \notin E^+
\]
\textit{(Không một thành viên nào có thể là tổ tiên của chính mình).}
\end{definition}

\begin{definition}[Ràng buộc C4.3: Tuổi của Cha/Mẹ khi sinh Con]
Với mọi cặp quan hệ $(p, c) \in \text{ParentChildRelation}$ liên kết với hồ sơ thành viên $m_{\text{parent}}$ và $m_{\text{child}}$:
\[
m_{\text{child}}.\text{birth\_date} - m_{\text{parent}}.\text{birth\_date} \ge 12\text{ (năm)}
\]
\end{definition}

\subsection{Bảng tổng hợp Ràng buộc Toàn vẹn và Cơ chế Cài đặt}
\label{subsec:2_4_5_tong_hop}

\begin{table}[htbp]
\centering
\small
\caption{Bảng tổng hợp Ràng buộc Toàn vẹn Mức Logic của Hệ thống FamilyConnect.}
\label{tab:integrity_matrix}
\begin{tabularx}{\textwidth}{p{2.6cm} X p{3.2cm} Y}
\toprule
\textbf{Mã Ràng buộc} & \textbf{Tên / Bối cảnh Nghiệp vụ} & \textbf{Loại Ràng buộc} & \textbf{Cơ chế Cài đặt (PostgreSQL)} \\
\midrule
\textbf{C1.1 -- C1.3} & Miền giá trị Giới tính, RSVP, Reaction & Domain Constraint & \texttt{CHECK} / \texttt{ENUM} Type \\
\textbf{C2.1 -- C2.3} & Duy nhất Email, Username, Khóa chính & Entity Constraint & \texttt{PRIMARY KEY}, \texttt{UNIQUE} \\
\textbf{C3.1 -- C3.2} & Tham chiếu Hồ sơ, Mối quan hệ Cha--Con & Referential Integrity & \texttt{FOREIGN KEY ON DELETE} \\
\textbf{C4.1} & Tính hợp lệ Ngày sinh -- Ngày mất & Semantic Constraint & \texttt{CHECK (death\_date >= birth\_date)} \\
\textbf{C4.2} & Không lặp tổ tiên (DAG Constraint) & Cross-Table Constraint & Trigger \texttt{BEFORE INSERT/UPDATE} \\
\textbf{C4.3} & Tuổi cha mẹ khi sinh con ($\ge 12$ tuổi) & Semantic Constraint & Trigger \texttt{BEFORE INSERT/UPDATE} \\
\bottomrule
\end{tabularx}
\end{table}

\subsection{Biểu diễn Đồ thị Quan hệ (Graph Data Model)}
\label{subsec:2_4_6_graph}

Ngoài mô hình quan hệ bảng chuẩn, hệ thống \textbf{FamilyConnect} bổ sung lớp biểu diễn dạng Đồ thị tương thích với Lớp dịch vụ AI để hỗ trợ tra cứu quan hệ huyết thống và trực quan hóa gia phả tương tác.

\begin{figure}[htbp]
\centering
\begin{tikzpicture}[
    node distance=1.8cm and 2.2cm,
    entity/.style={rectangle, draw=blue, fill=blue!5, thick, minimum width=2.6cm, minimum height=1.1cm, align=center, font=\small},
    rel/.style={diamond, draw=red!80!black, fill=red!5, thick, aspect=1.8, inner sep=2pt, align=center, font=\small},
    line/.style={draw, thick, -{Stealth[length=2.5mm]}}
]
    \node[entity] (member) {\textbf{MemberProfile}\\(\underline{member\_id})};
    \node[rel, above=1.5cm of member] (marriage) {Marriage};
    \node[entity, left=2.2cm of member] (user) {\textbf{Users}\\(\underline{user\_id})};
    \node[entity, right=2.2cm of member] (family) {\textbf{FamilyTree}\\(\underline{family\_id})};
    \node[entity, below left=2.0cm and 0.8cm of member] (post) {\textbf{Post}\\(\underline{post\_id})};
    \node[entity, below right=2.0cm and 0.8cm of member] (event) {\textbf{Event}\\(\underline{event\_id})};
    \node[rel, below=1.8cm of member] (parent) {Parent-Child\\(Recursive)};

    \draw[line] (user) -- node[above, font=\footnotesize] {1:1} (member);
    \draw[line] (member) -- node[above, font=\footnotesize] {N:1} (family);
    \draw[line] (member) -- node[left, font=\footnotesize] {1:N} (post);
    \draw[line] (member) -- node[right, font=\footnotesize] {1:N} (event);
    \draw[line] (member.north) -- (marriage.south);
    \draw[line] (member.south) -- (parent.north);
\end{tikzpicture}
\caption{Sơ đồ các mối liên kết cốt lõi của hệ thống FamilyConnect (Graph Data Model).}
\label{fig:family_logic_graph}
\end{figure}



% Chương 3: Thiết kế Vật lý (Physical Design)
% !TEX root = ../../main.tex
\chapter{Thiết kế Vật lý (Physical Design)}
\label{chap:chuong3_vat_ly}

Giai đoạn thiết kế vật lý hiện thực hóa lược đồ quan hệ logic lên một hệ quản trị cơ sở dữ liệu cụ thể (MySQL 8.0+), xác lập từ điển dữ liệu chi tiết, cài đặt các chỉ mục tối ưu hóa truy vấn, xây dựng các thủ tục/trigger nâng cao và hoạch định chiến lược an toàn dữ liệu.

% !TEX root = ../../main.tex
\section{Lựa chọn Hệ quản trị CSDL và công nghệ}
\label{sec:3_1_lua_chon_dbms}

\subsection{Lựa chọn MySQL và lưu trữ quan hệ (ACID, CTE)}
% Nội dung lý do lựa chọn MySQL

\subsection{Công nghệ mở rộng (Vector DB, S3 / File Storage)}
% Nội dung công nghệ mở rộng

% !TEX root = ../../main.tex
\section{Từ điển dữ liệu chi tiết (Data Dictionary)}
\label{sec:3_2_tu_dien_du_lieu}

% Bảng mô tả chi tiết từng cột: Tên trường, Kiểu dữ liệu vật lý (UUID, VARCHAR, TIMESTAMP, JSON,...), Null?, Default, Mô tả ý nghĩa.

% !TEX root = ../../main.tex
\section{Chiến lược đánh chỉ mục (Indexing Strategy)}
\label{sec:3_3_chien_luoc_index}

\subsection{Chỉ mục Khóa chính và Khóa ngoại (B-Tree Indexing)}
% Nội dung đánh index cho PK, FK tối ưu JOIN

\subsection{Chỉ mục tối ưu hóa Truy vấn Phả hệ và Tìm kiếm}
% Nội dung chỉ mục cho clan_id, branch_id, created_at, họ tên, số điện thoại, duyệt cây

% !TEX root = ../../main.tex
\section{Thiết kế các đối tượng nâng cao (Views, Stored Procedures, Functions, Triggers)}
\label{sec:3_4_doi_tuong_nang_cao}

\subsection{Trigger kiểm soát toàn vẹn Cây phả hệ (Anti-Loop Trigger)}
% Nội dung Trigger chống vòng lặp tổ tiên

\subsection{Stored Procedure duyệt cây phả hệ đệ quy (Recursive CTE Traversal)}
% Nội dung Stored Procedure / Function lấy toàn bộ cây phả hệ hoặc tìm đường đi quan hệ

\subsection{View tổng hợp Thống kê và Báo cáo}
% Nội dung các View báo cáo

% !TEX root = ../../main.tex
\section{Kịch bản DDL tạo bảng (Physical Database Scripts)}
\label{sec:3_5_ddl_scripts}

% Các đoạn mã SQL CREATE TABLE, ALTER TABLE ADD CONSTRAINT tiêu biểu

% !TEX root = ../../main.tex
\section{An toàn dữ liệu \& Chiến lược sao lưu}
\label{sec:3_6_an_toan_sao_luu}

\subsection{Phân quyền truy cập tầng Cơ sở dữ liệu (Roles / Grants)}
% Nội dung phân quyền truy cập CSDL

\subsection{Chiến lược Ghi nhật ký kiểm toán (Audit Logging)}
% Nội dung Audit Logging

\subsection{Chiến lược Sao lưu và Phục hồi (Backup \& Disaster Recovery)}
% Nội dung sao lưu và phục hồi định kỳ



% Đánh giá và Kết luận (Không đánh số chương)
% !TEX root = ../main.tex
\chapter*{ĐÁNH GIÁ VÀ KẾT LUẬN}
\addcontentsline{toc}{chapter}{ĐÁNH GIÁ VÀ KẾT LUẬN}
\label{chap:danh_gia_ket_luan}

\section*{1. Kết quả đạt được}
\addcontentsline{toc}{section}{1. Kết quả đạt được}
% Nội dung kết quả đạt được

\section*{2. Hạn chế còn tồn tại}
\addcontentsline{toc}{section}{2. Hạn chế còn tồn tại}
% Nội dung hạn chế còn tồn tại

\section*{3. Hướng phát triển tiếp theo}
\addcontentsline{toc}{section}{3. Hướng phát triển tiếp theo}
% Nội dung hướng phát triển tiếp theo


% ------------------------------------------------------------------------------
% TÀI LIỆU THAM KHẢO
% ------------------------------------------------------------------------------
\begin{thebibliography}{99}
\bibitem{chen1976}
Peter Pin-Shan Chen.
\newblock The Entity-Relationship Model---Toward a Unified View of Data.
\newblock {\em ACM Transactions on Database Systems (TODS)}, 1(1):9--36, 1976.

\bibitem{armstrong1974}
William Ward Armstrong.
\newblock Dependency structures of data base relationships.
\newblock {\em Information Processing 74}, pages 580--583. North-Holland, 1974.

\bibitem{elmasri2015}
Ramez Elmasri and Shamkant B. Navathe.
\newblock {\em Fundamentals of Database Systems}.
\newblock Pearson, 7th edition, 2015.

\bibitem{silberschatz2019}
Abraham Silberschatz, Henry F. Korth, and S. Sudarshan.
\newblock {\em Database System Concepts}.
\newblock McGraw-Hill, 7th edition, 2019.
\end{thebibliography}

\end{document}
